\documentclass{article}
\newcommand{\ManualName}{proccontrol}

\usepackage{pgf}
\pgfrealjobname{proccontrol}

\usepackage{amssymb}


\usepackage{tikz}
\usetikzlibrary{arrows,positioning,fit,shapes,decorations,decorations.pathmorphing,decorations.pathreplacing,calc,patterns,scopes,matrix,backgrounds}
\usepackage{booktabs}
%% tabularx lets a table size its own columns: l takes each column's
%% natural width and X absorbs the rest, so the table fills \textwidth
%% without any hand-computed widths to go stale.
\usepackage{tabularx}
\usepackage{xcolor}
\usepackage[hyperindex,breaklinks]{hyperref}
\usepackage{helvet}
\usepackage{enumitem}
\usepackage[letterpaper,margin=1in]{geometry}
\usepackage{listings}
\usepackage{verbatim}
\usepackage[T1]{fontenc}

\setlength{\parindent}{0.0in}
\setlength{\parskip}{0.1in}

\newenvironment{apient}{\small\verbatim}{\endverbatim}

\newcommand{\apidesc}[1]{%
{\addtolength{\leftskip}{4em}%
#1\par\medskip}
}

\newcommand{\declaredin}[1]{%
\textbf{Declared in:} \code{#1}
}

\newcommand{\code}[1]{{\small\texttt{#1}}}

%% Comments are meant to recede without becoming unreadable; a muted blue
%% of roughly the lightness the old lightgray had.
\definecolor{commentblue}{rgb}{0.40,0.55,0.75}

%% One style for every code listing in every manual.  Keep it here: an
%% \lstset outside a group is global and lasts to the end of the document,
%% so a stray one in an input file silently restyles every later listing.
%% A manual adds only its own emph list; everything else belongs here.
%%
%% breakatwhitespace keeps a wrap from landing inside an identifier, and
%% postbreak marks the continuation so a wrapped line cannot be mistaken
%% for one the author wrote that way.
\lstset{
  language={[GNU]C++},
  basicstyle=\ttfamily\small,
  keywordstyle=\ttfamily\color{purple},
  commentstyle=\ttfamily\color{commentblue},
  stringstyle=\ttfamily\color{orange},
  emphstyle=\color{blue},
  frame=leftline,
  numbers=left,
  numberstyle=\tiny,
  stepnumber=3,
  numbersep=5pt,
  showstringspaces=false,
  breaklines=true,
  breakatwhitespace=true,
  postbreak=\mbox{\textcolor{red}{$\hookrightarrow$}\space},
}


\begin{document}
%% article.cls sets the page counter to 1 both on entering and on leaving
%% titlepage, so the title page and the page after it are both page 1 and
%% hyperref emits a duplicate page.1 destination.  Give the title page no
%% page anchor; nothing links to it.
\hypersetup{pageanchor=false}
\begin{titlepage}
\beginpgfgraphicnamed{titlepage}
\begin{tikzpicture}[remember picture, overlay]
    \path
        (current page.north west) coordinate (origin)
        (current page.north) coordinate (topcenter);

    % Header
    \node [font=\sffamily] (pppt) at ($(topcenter) - (0,1.0in)$)
        {\fontsize{24}{36}\selectfont Paradyn Parallel Performance Tools};

    % Document Title
% older versions of pgf have a bug for matrices in overlay mode;
% have to specify positions manually
%    \matrix (Title) [%
%        matrix of nodes,%
%        nodes={font=\sffamily,right},%
%        matrix anchor=west,%
%        row sep=12pt
%        ] at ($(origin)+(0.75in,-3.0in)$)
%    {
%        \fontsize{48}{56}\selectfont ParseAPI \\
%        \fontsize{44}{56}\selectfont Programmer's Guide \\
%    };

    \node [anchor=west,font=\sffamily] (title1) at ($(origin)+(0.75in,-3.0in)$)
        {\fontsize{48}{56}\selectfont \ManualName};
    \node [anchor=west,font=\sffamily] (title2) at ($(title1.west)+(0in,-56pt)$)
        {\fontsize{48}{56}\selectfont Programmer's Guide};

    % Release information
%    \matrix (Releaseinfo) [%
%        matrix of nodes,%
%        nodes={font=\sffamily,right},%
%        matrix anchor=west,%
%        row sep=8pt
%        ] at ($(origin)+(0.75in,-5.0in)$)
%    {
%        %\fontsize{24}{32}\selectfont Release 0.1 \\
%        \fontsize{24}{32}\selectfont Beta Release \\
%        \fontsize{24}{32}\selectfont Oct 2010 \\
%    };

    \node [anchor=west,font=\sffamily] (rel1) at ($(origin)+(0.75in,-5.0in)$)
        {\fontsize{24}{32}\selectfont 13.0 Release};
    \node [anchor=west,font=\sffamily] (rel2) at ($(rel1.west)+(0in,-32pt)$)
        {\fontsize{24}{32}\selectfont February 2024};

    % Contact information
%    \matrix (UWaddress) [%
%        matrix of nodes,%
%        nodes={font=\sffamily\large,right},%
%        matrix anchor=north west
%        ] at ($(origin)+(0.75in,-7in)$)
%    {
%        Computer Science Department \\
%        University of Wisconsin--Madison \\
%        Madison, WI 53706 \\
%    };

    \node [anchor=west,font=\sffamily\large] (uw1) at ($(origin)+(0.75in,-7.0in)$)
        {Computer Sciences Department};
    \node [anchor=west,font=\sffamily\large] (uw2) at ($(uw1.west)+(0in,-20pt)$)
        {University of Wisconsin--Madison};
    \node [anchor=west,font=\sffamily\large] (uw3) at ($(uw2.west)+(0in,-20pt)$)
        {Madison, WI 53706};


%    \matrix (UMDaddress) [%
%        matrix of nodes,%
%        nodes={font=\sffamily\large,right},%
%        matrix anchor=north west,
%        below=1em of UWaddress.south west
%        ]
%    {
%        Computer Science Department \\
%        University of Maryland \\
%        College Park, MD 20742 \\
%    };

    \node [anchor=west,font=\sffamily\large] (umd1) at ($(uw3.south west)+(0in,-2.5em)$)
        {Computer Science Department};
    \node [anchor=west,font=\sffamily\large] (umd2) at ($(umd1.west)+(0in,-20pt)$)
        {University of Maryland};
    \node [anchor=west,font=\sffamily\large] (umd3) at ($(umd2.west)+(0in,-20pt)$)
        {College Park, MD 20742};

%    \matrix (Emails) [%
%        matrix of nodes,%
%        nodes={font=\sffamily,right},%
%        matrix anchor=north west,%
%        below=1em of UMDaddress.south west,%
%        anchor=base
%        ]
%    {
%        Email & \code{bugs@dyninst.org} \\
%        Web & \code{www.dyninst.org} \\
%    };

    \node [anchor=west,font=\sffamily] (email1) at ($(umd3.south west)+(-0.5em,-2.5em)$)
        %{Email \code{bugs@dyninst.org}};
        {\begin{tabular}{ll}%
         Email & \texttt{dyninst-team@g-groups.wisc.edu} \\
         Web & \texttt{https://github.com/dyninst/dyninst} \\
        \end{tabular}};


    % Logo
    \path
        node (logo) at ($(origin)+(4.0in,-7.0in)$) [%
            anchor=north west]
        {%
            \includegraphics[width=3.25in]{paradyn_logo}
        };


\end{tikzpicture}
\endpgfgraphicnamed
\end{titlepage}
\hypersetup{pageanchor=true}


\lstset{
  %% Highlight this manual's own type names like keywords; the style
  %% itself lives in common/manual-latex/manual_commands.tex.
  emph={
    Process,
    Event,
    EventType
  },
}

\tableofcontents
\clearpage

\setcounter{secnumdepth}{5}

\section{Introduction}
  This document describes ProcControlAPI, an API and library for controlling
processes. ProcControlAPI runs as part of a controller process
and manages one or more target processes. It allows the controller process
to perform operations on target
processes, such as writing to memory, stopping and running threads, or
receiving notification when certain events occur. ProcControlAPI presents
these operations through a platform-independent API and high-level
abstractions. Users can describe what they want ProcControlAPI to do, and
ProcControlAPI handles the details.

An example use for ProcControlAPI would be as the underlying mechanism for a
debugger. A user writing a debugger could provide their own user interface and
debugging strategies, while using ProcControlAPI to perform operations such as
creating processes, running threads, and handling breakpoints.

The interface for ProcControlAPI can be generally divided into two parts: an
interface for managing a process (e.g., reading and writing to target process
memory, stopping and running threads), and an interface for monitoring a target
process for certain events (e.g., watching the target process for fork or
thread creation events). The manager interface uses a set of C++ objects to
represent a target process and its threads, libraries, registers and other
interesting aspects. Operations performed on these C++ objects in the
controller process are translated into corresponding operations on the target
process. The event interface uses a callback system to notify the
ProcControlAPI user of interesting events in the target process.

\subsection{Simple Example}

As an example, consider the code in Listing~\ref{lst:thread-create} that creates a target process and
prints a message whenever that target process creates a new thread. Details
on the API function used in this example can be found in latter sections of
this manual, but we will provide a high level description of the operations
here. Note that proper error handling and checking have been left out for
brevity.

\lstinputlisting[label=lst:thread-create, caption={Printing a message when the target process creates a thread}]{examples/ThreadCreate.C}

We ask ProcControlAPI to create a new Process using the
given arguments. ProcControlAPI will spawn a new target process
 and leave it in a stopped state to prevent it from executing.

\begin{enumerate}
\item After creating the new target process we register a callback
function. We ask ProcControlAPI to call our
function, on\_thread\_create, when an event of type EventType::ThreadCreate
occurs in the target process.

\item The on\_thread\_create function takes a pointer to an object of type
Event and returns a Process::cb\_ret\_t. The Event describes the target
process event that triggered
this callback. In this case, it provides information about the new thread in
the target process. It is worth noting that Event::const\_ptr is a not a
regular pointer, but a reference counted shared pointer. This means that we
do not have to be concerned with cleaning the Event- it
will be automatically cleaned when the last reference disappears. The
Process::cb\_ret\_t describes what action should be taken on the process in
response to this event, which is described in more detail in
Section~\ref{subsec:process}.

\item The Event class has several child classes, one of which is EventNewThread.
We start by casting the Event into an EventNewThread and then extract
information about the new thread from the EventNewThread.

\item In step 6, we've finished handling the new thread
event and need to tell ProcControlAPI what to do in response to this event.
For example, we could choose to stop the process from further execution by
returning a value of Process::cbProcStop. Instead, we choose let
ProcControlAPI take its default action for an EventNewThread by returning
Process::cbDefault, which is to continue the process and its new thread (which
were both stopped before delivery of the callback).

\item The registering of our callback in step 3 did not actually trigger any
calls to the callback function- the target process was
created in a stopped state and has not yet been able to create any threads.
We tell ProcControlAPI to continue the target process in this step, which
allows it to execute and possibly start generating new events.

\item In this step we wait for the target process to finish executing and
terminate. Calling Process::handleEvents
blocks the controller process until an event occurs, allowing us to wait for
events without needing to spin the controller process on the CPU.
\end{enumerate}


\section{Abstractions}
  This section focuses on some of the more important concepts in ProcControlAPI
and gives a high level overview before the detailed API is presented,
beginning in Section~\ref{sec:processes}.

\subsection{Processes and Threads}
There are two central classes to ProcControlAPI, Process and Thread.
Each class respectively represents a single
target process or thread running on the system. By performing operations on
the Process and Thread objects, a ProcControlAPI user is able to control the
target process and its threads.

Each Process is guaranteed to have at least one Thread associated with it. A
multi-threaded process may have a Process object with more than one Thread.
Each process has an address space associated with it, which can be written or
read through the Process object. Each thread has a set of registers
associated with it, which can be access through the Thread object.

At any one time a Thread will be in either a stopped state or a running
state. A thread in a
stopped state has had its execution paused by
ProcControlAPI- the OS will not schedule the thread to run.
A thread in a running state is allowed to execute as normal. A thread in a
running state may block for other reasons, e.g. blocking on IO calls, but this
does not affect ProcControlAPI's view of the thread state.
A thread is only in the stopped state if ProcControlAPI has explicitly
stopped it.

A Process object is not considered to have a stopped or running
state- only its Thread objects are stopped or running. A
stop operation on a Process triggers a stop operation on each of its Threads,
and similarly a continue operation on a Process triggers continue operations on
each Thread.

\subsection{Callbacks}
In addition to controlling a target process through the Process and Thread
objects, a ProcControlAPI user can also receive notification of events that
happen in that process. Examples of these events would be a new thread being
created, a breakpoint being executed, or a process exiting.

The ProcControlAPI user receives notice of events through a callback system.
The user can register callback function that will be called by ProcControlAPI
whenever a particular type of event occurs. Details about the event are
passed to the callback function via an Event object.

\subsubsection{Events}
\label{subsubsec:events}
Each event can be broken up into
an EventType object and an Event
object. The EventType describes a type of event that can happen, and Event
describes a specific instance of an event happening. Each Event will have one
and only one EventType.

Each EventType has two primary fields: its time and its code. The code field
of describes what type of event occurred, e.g. EventType::Exit represents a
target process exiting. The time field of an EventType represents whether the
EventType is happening before or after will have code and will have a value of
EventType::Pre, EventType::Post, or EventType::None.

For example, an EventType with time and code of EventType::Pre and
EventType::Exit will occur just before a target process exits, and a code of
EventType::Exec with a time of EventType::Post will occur after an exec system
call occurs. In this document we will abbreviate EventTypes such as these as
pre-exit and post-exec. Some EventTypes do not have a time associated with
them, for example EventType::Breakpoint does not have an associated time and
thus has a time value of EventType::none.

An Event represents an instance of an EventType occurring. In addition to an
EventType, each Event also has pointer to the Process and Thread that it
occurred on. Certain events may also have event specific information
associated with them, which is represented in a sub-class of Event. Each
EventType is associated with a specific sub-class of Event.

For example, EventType::Library is used to signify a shared library being loaded
into the target process. When an EventType::Library occurs ProcControlAPI
will deliver an object of type EventLibrary, which is a subclass of Event, to
any registered callback functions. In addition to the information inherited
from Event, the EventLibrary will contain extra information about the library
that was loaded into the target process.

Table~\ref{tab:event-subclass} shows the Event subclass that is used for each EventType. Not all
EventTypes are available on every platform- a checkmark
under the specific OS column means that the EventType is available on that OS.

\begin{table}[!ht]
  \centering
  \begin{tabular}{|l|l|l|l|l|}
  \hline
    EventType & Event Subclass & Linux & FreeBSD & Windows \\ \hline
    Stop & EventStop & \checkmark & ~ & ~ \\ \hline
    Breakpoint & EventBreakpoint & \checkmark & \checkmark & \checkmark \\ \hline
    Signal & EventSignal & \checkmark & \checkmark & \checkmark \\ \hline
    UserThreadCreate & EventNewUserThread & \checkmark & \checkmark & ~ \\ \hline
    LWPCreate & EventNewLWP & \checkmark & ~ & \checkmark \\ \hline
    Pre-UserThreadDestroy & EventUserThreadDestroy & \checkmark & \checkmark & ~ \\ \hline
    Post-UserThreadDestroy & EventUserThreadDestroy & \checkmark & \checkmark & ~ \\ \hline
    Pre-LWPDestroy & EventLWPDestroy & \checkmark & ~ & \checkmark \\ \hline
    Post-LWPDestroy & EventLWPDestroy & \checkmark & ~ & ~ \\ \hline
    Pre-Fork & EventFork & ~ & ~ & ~ \\ \hline
    Post-Fork & EventFork & \checkmark & ~ & ~ \\ \hline
    Pre-Exec & EventExec & ~ & ~ & ~ \\ \hline
    Post-Exec & EventExec & \checkmark & \checkmark & ~ \\ \hline
    RPC & EventRPC & \checkmark & \checkmark & \checkmark \\ \hline
    SingleStep & EventSingleStep & \checkmark & \checkmark & \checkmark \\ \hline
    Breakpoint & EventBreakpoint & \checkmark & \checkmark & \checkmark \\ \hline
    Library & EventLibrary & \checkmark & \checkmark & \checkmark \\ \hline
    Pre-Exit & EventExit & \checkmark & ~ & ~ \\ \hline
    Post-Exit & EventExit & \checkmark & \checkmark & \checkmark \\ \hline
    Crash & EventCrash & \checkmark & \checkmark & \checkmark \\ \hline
    ForceTerminate & EventForceTerminate & \checkmark & \checkmark & \checkmark \\ \hline
  \end{tabular}
  \caption{Event subclass used for each EventType}
  \label{tab:event-subclass}
\end{table}


\subsubsection{Callback Functions}
Events are delivered via a
callback function. A ProcControlAPI user can register callback functions for
an EventType using the Process::registerEventCallback function. All callback
functions must be declared using the signature:

\begin{lstlisting}
Process::cb_ret_t callback_func_name(Event::ptr)
\end{lstlisting}

In order to prevent a class of race conditions, ProcControlAPI does not allow a
callback function to perform any operation that would require another callback
to be recursively delivered. At most one callback function can be running at
a time.

To enforce this, the event that is passed to a callback function contains only
const pointers to the triggering Process and Thread objects. Any member
function that could trigger callbacks is not marked const, thus triggering a
compilation error if they are called on an object passed to a callback. If
the ProcControlAPI user uses const\_cast or global variables to get around the
const restriction it will result in a runtime error. API functions that
cannot be used from a callback are mentioned in the API entries.

Operations such as Process::stopProc, Process::continueProc, Thread::stopThread, and
Thread::continueThread are not safe to
call from a callback function, but it would still be useful to perform these
operations. ProcControlAPI allows the user to use the return value from a
callback function to specify whether process or thread that triggered the event
should be stopped or continued. More details on this can be found in the
\code{Process::cb\_ret\_t} section of the API reference.

\subsubsection{Callback Delivery}
\label{subsubsec:callback-delivery}
When ProcControlAPI needs to
deliver a callback it must first gain control of a user visible thread in the
controller process. This thread will be used to invoke the callback
function. ProcControlAPI does not use its
internal threads for delivering callbacks, as this would expose the
ProcControlAPI user to race conditions.

Unfortunately, the user thread is not always accessible to ProcControlAPI when
it needs to invoke a callback function. For example, the user visible thread
may be performing network IO or waiting for input from a GUI when an event
occurs.

ProcControlAPI uses a notification system built around the EventNotify class to
alert the ProcControlAPI user that a callback is ready to be delivered. Once
the user is notified then they can call the Process::handleEvents function,
under which ProcControlAPI will invoke any pending callback functions.

The EventNotify class has two mechanisms for
notifying the ProcControlAPI user that a callback is pending: writing to a file
descriptor and a light-weight callback function. The
EventNotify::getFD function returns a file
descriptor that will have a byte written to it when a callback is ready. This
file descriptor can be added to a select or poll to block a thread that handles
ProcControlAPI events. Alternatively, the ProcControlAPI user can register a
light-weight callback that is invoked when a callback is ready. This
light-weight callback provides no information about the Event and may occur on
another thread or from a signal handler- the ProcControlAPI
user is encouraged to keep this callback minimal.

It is important for a user to respond promptly to a callback notification. A
target process may remain blocked while a notification is pending. If a
target process is generating many events that need callbacks, a long delay in
notification could have a significant performance impact.

Once the ProcControlAPI user knows that a callback is ready to be delivered they
can call Process::handleEvents, which will
invoke all callback functions. Alternatively, if the ProcControlAPI user does
not need to handle events outside of ProcControlAPI, they can continue to block
in Process::handleEvents without going through the notification system.

\subsection{Inferior Remote Procedure Call}
An iRPC is a mechanism for executing code in a target process. Despite the
name, an iRPC does not necessarily have to involve a procedure
call- any piece of code can be executed.

A ProcControlAPI user can invoke an iRPC by providing ProcControlAPI with a
buffer of machine code and specifying a Process or Thread on which to run the
machine code. ProcControlAPI will insert the machine code into the address
space, save the register set, run the machine code, and then remove the machine
code after execution completes. When the iRPC completes (but before the
registers and memory are cleaned) ProcControlAPI will deliver an
EventIRPC to any registered callback function. The
ProcControlAPI user may use this callback to collect any results from the
registers or memory used by the iRPC.

Note that ProcControlAPI will preserve the registers of the thread running the
iRPC, and it will preserve the memory used by the machine code. Other memory
or system state changed by the iRPC may remain visible to the target process
after the iRPC completes.

The machine code for each iRPC must contain at least one trap instruction (e.g.,
a 0xCC instruction on x86 family or a 0x7D821008 instruction on the PPC
family). ProcControlAPI will stop executing the iRPC upon invocation of the
trap. Note that the trap instruction must fall within the original machine
code for the iRPC. If the iRPC calls or jumps to another piece of code that
executes a trap instruction then ProcControlAPI will not treat it as the end of
the iRPC.

Before an iRPC can be run it must be posted to a process or thread using the
Process::postIRPC or
Thread::postIRPC API functions. The
Process::postIRPC function will select a thread to post the iRPC to.
Multiple iRPCs can be posted to the same thread, but only one iRPC will run
at a time- subsequent iRPCs will be queued and run after
the preceding iRPC completes. If multiple iRPCs are posted to different
threads in a multi-threaded process, then they may run in parallel.

An iRPC can be posted to a stopped or running thread. If posted to a stopped
thread, then the iRPC will run when the thread is continued. If posted to a
running thread, then the iRPC will run immediately or, if posted from a
callback function, when the callback function completes.

An iRPC may be blocking or non-blocking. If a blocking iRPC is posted to any
Process, then calls to Process::handleEvents will block until the iRPC is completed.


\section{Processes}
\label{sec:processes}
  ProcControlAPI offers several facilities for interacting with processes in a
  platform-independent manner.

  \subsection{Class Process}
  \label{subsec:process}
    \declaredin{PCProcess.h}

The Process class is the primary handle for operating on a single target process.
Process objects may be created by calls to the static functions
Process::createProcess or Process::attachProcess, or in response to
certain types of events (e.g, fork on UNIX systems).

The static functions of the Process class serve as a central location for
performing general ProcControlAPI operations, such as handleEvents and
registerEventCallback when dealing with callbacks.

\begin{apient}
Process::ptr;
Process::const_ptr;
\end{apient}

\apidesc{
The Process::ptr and Process::const\_ptr respectively represent a pointer and a
const pointer to a Process object. Both pointer types are reference counted
and will cause the underlying Process object to be cleaned when there are no
more references. ProcControlAPI will maintain internal references to any
Process it actively controls, relinquishing those references when the process
either exits or is detached.
}

\begin{apient}
enum Process::cb_action_t {
  cbDefault,
  cbThreadContinue,
  cbThreadStop,
  cbProcContinue,
  cbProcStop
};

struct Process::cb_ret_t {
  cb_ret_t(cb_action_t p);
  cb_ret_t(cb_action_t p, cb_action_t c);
  cb_action_t parent;
  cb_action_t child;
};
\end{apient}

\apidesc{
The cb\_ret\_t struct is used as the return type for callback functions registered
through Process::registerEventCallback(). A callback function can specify
whether the thread or process associated with its event should be stopped or
continued by respectively returning cbThreadContinue, cbThreadStop,
cbProcContinue, or cbProcStop. The cbDefault return value returns a Process
and Thread to the original state before the event occurred.
Some events, such as process spawn or thread create involve two processes or
threads. In this case the ProcControlAPI user can specify a cb\_action\_t
value for both the parent and child using the two parameter constructor for
cb\_ret\_t.
}

\begin{apient}
typedef Process::cb_ret_t(*cb_func_t)(Event::const_ptr);
\end{apient}

\apidesc{
The cb\_func\_t type is a function pointer type for functions that can handle
event callbacks. The callback function gets an Event::const\_ptr as input,
which points to the Event that triggered the callback. The cb\_func\_t
function should return a cb\_ret\_t describing what to do with the process
after handling the event.
}

\begin{apient}
enum OSType{
  OSNone,
  Linux,
  FreeBSD,
  Windows
};
\end{apient}

\apidesc{
A value from this enum is returned from Process::getOS and signifies the current
OS on which the target process is running.
}

\begin{apient}
typedef std::pair<Dyninst::Address, Dyninst::Address> MemoryRegion;
\end{apient}

\apidesc{
The MemoryRegion type represents a region of allocated memory, and the first
part of the pair is the start address, the second, the end.
}

\begin{apient}
static Process::ptr createProcess(std::string executable,
                                  const std::vector<std::string>& argv,
                                  const std::vector<std::string>& envp = emptyEnvp,
                                  const std::map<int, int>& fds = emptyFDs);
\end{apient}

\apidesc{
This function creates a new process by launching an executable file named by
executable with the arguments specified by argv, the environment specified in
envp, and it returns a pointer to the new Process object upon success. The
new process will be created with its initial thread in the stopped state.
It is an error to call this function from a callback.
If the fds map is not empty, then the new process will be created with the file
descriptors from the fds' first elements dup2 mapped to
the file descriptors in fds' second elements.
If envp is empty, the environment will be inherited from the calling process.
ProcControlAPI may deliver callbacks when this function is called.
This function returns Process::ptr() on error, and a subsequent call to
getLastError returns details on the error.
}

\begin{apient}
static Process::ptr attachProcess(Dyninst::PID pid, std::string executable = "");
\end{apient}

\apidesc{
This function creates a new Process object by attaching to the PID specified by
pid. The new Process object will be returned from this function upon success.
The executable argument is optional, and can be used to assist ProcControlAPI
in finding the process' executable on operating systems
where this cannot be easily determined (currently on AIX). The new process
will be returned with all of its threads in the stopped state.
It is an error to call this function from a callback.
ProcControlAPI may deliver callbacks when this function is called.
This function return Process::ptr() on error, and a subsequent call to
getLastError returns details on the error.
}

\begin{apient}
static bool handleEvents(bool block);
\end{apient}

\apidesc{
This function causes ProcControlAPI to handle any pending debug events and
deliver callbacks. When an event requires a
callback ProcControlAPI needs control of the main thread in order to deliver
the callback. This function gives control of the main thread to
ProcControlAPI for callback delivery. A user can know when to call
handleEvents by using the EventNotify interface.
If the block parameter is true, then handleEvents will block until at least one
debug event has been handled. If block is false then handleEvents returns
immediately if no events are ready to be handled.
This function returns true if it handled at least one event and false otherwise.
It is an error to call this function from a callback.
}

\begin{apient}
static bool registerEventCallback(EventType evt, cb_func_t cbfunc);
\end{apient}

\apidesc{
This function registers a new callback function with ProcControlAPI. Upon
receiving an event with type evt, ProcControlAPI will deliver a callback with
that event to the cbfunc function. Multiple functions can be registered to
receive callbacks for a single EventType, and a single function can be
registered with multiple EventTypes.
If multiple callback functions are registered with a single EventType, then it
is undefined what order those callback functions will be invoked in. In this
case the cb\_ret\_t result of the last callback function called will be used to
determine what stop or continue operations should be performed on the process.
If a single callback function is registered for the same EventType multiple
times, then ProcControlAPI will only invoke one call to the callback function
for each instance of the EventType.
This function returns true on success and false on error. Upon an error a
subsequent call to getLastError returns details on the error.
}

\begin{apient}
static bool removeEventCallback(EventType evt, cb_func_t cbfunc);
\end{apient}

\apidesc{
This function un-registers a callback that was registered with
registerEventCallback. After a successful call to this function the callback
function cbfunc will stop being called for events with EventType evt. Other
callback functions registered for evt will not be affected. Other instances
of cbfunc registered for different EventTypes will not be affected.
This function returns true if a callback was successfully removed and false
otherwise. Upon an error a subsequent call to getLastError returns details on
the error.
}

\begin{apient}
static bool removeEventCallback(EventType evt);
\end{apient}

\apidesc{
This function unregisters all callback functions associated with the EventType
evt. After a successful call to this function ProcControlAPI will stop
delivering callbacks for evt until a new callback function is registered.
This function returns true if a callback was successfully removed and false
otherwise. Upon an error a subsequent call to getLastError returns details on
the error.
}

\begin{apient}
static bool removeEventCallback(cb_func_t func);
\end{apient}

\apidesc{
This function unregisters all instances of callback function func from any
callback with any EventType.
This function returns true if a callback was successfully removed and false
otherwise. Upon an error a subsequent call to getLastError returns details on
the error.
}

\begin{apient}
Dyninst::PID getPid() const;
\end{apient}

\apidesc{
This function returns an OS handle referencing the process. On UNIX systems
this is the pid of the process.
}

\begin{apient}
Dyninst::Architecture getArchitecture() const;
\end{apient}

\apidesc{
This function returns an enum that describes the architecture of the target
process.
}

\begin{apient}
Dyninst::OSType getOS() const;
\end{apient}

\apidesc{
This function returns an enum that describes the OS of the target process. See
the beginning of this section for the definition of Dyninst::OSType.
}

\begin{apient}
bool supportsLWPEvents() const;
\end{apient}

\apidesc{
This function returns true if the target process can throw LWP create and
destroy events and false otherwise.
}

\begin{apient}
bool supportsUserThreadEvents() const;
\end{apient}

\apidesc{
This function returns true if the target process can throw user thread create
and destroy events and false otherwise.
}

\begin{apient}
bool supportsFork() const;
\end{apient}

\apidesc{
This function returns true if the fork system call is supported in the target
process and false otherwise.
}

\begin{apient}
bool supportsExec() const;
\end{apient}

\apidesc{
This function returns true if the exec system call is supported in the target
process and false otherwise.
}

\begin{apient}
bool isTerminated() const;
\end{apient}

\apidesc{
This function returns true if the target process has terminated (either via a
crash or normal exit) or if the ProcControlAPI has detached from the target
process. It returns false otherwise.
}

\begin{apient}
bool isExited() const;
\end{apient}

\apidesc{
This function returns true of the target process exited via a normal exit (e.g,
calling the exit function or returning from main). It returns false
otherwise.
}

\begin{apient}
int getExitCode() const;
\end{apient}

\apidesc{
If a target process exited normally then this function returns its exit code.
The return result of this function is undefined if the
Process' isExited function returns false.
}

\begin{apient}
bool isCrashed() const;
\end{apient}

\apidesc{
This function returns true if the target process exited because of a crash. It
returns false otherwise.
}

\begin{apient}
int getCrashSignal() const;
\end{apient}

\apidesc{
If a target process exited because of a crash, then this function returns the
signal that caused the target process to crash. The return result of this
function is undefined if the Process' isCrashed function
returns false.
}

\begin{apient}
bool hasStoppedThread() const;
\end{apient}

\apidesc{
This function returns true if the target process has at least one thread in the
stopped state. It returns false otherwise or if an error occurs. In the
event of an error a call to getLastError returns details on the error.
}

\begin{apient}
bool hasRunningThread() const;
\end{apient}

\apidesc{
This function returns true if the target process has at least one thread in the
running state. It returns false otherwise or if an error occurs. In the
event of an error a call to getLastError returns details on the error.
}

\begin{apient}
bool allThreadsStopped() const;
\end{apient}

\apidesc{
This function returns true if all threads in the target process are in the
stopped state. It returns false otherwise or if an error occurs. In the
event of an error a call to getLastError returns details on the error.
}

\begin{apient}
bool allThreadsRunning() const;
\end{apient}

\apidesc{
This function returns true if all threads in the target process are in the
running state. It returns false otherwise or if an error occurs. In the
event of an error a call to getLastError returns details on the error.
}

\begin{apient}
bool allThreadsRunningWhenAttached() const;
\end{apient}

\apidesc{
This function returns true if all threads were running when the controller
process attached to this process. It returns false if any threads were stopped.
If the target process was created instead of attached, this function returns
true.
}

\begin{apient}
bool continueProc();
\end{apient}

\apidesc{
This function will move all threads in the target process into the running
state. This function returns true if at least one thread was continued as
part of the call, and false otherwise.
It is an error to call this function from a callback.
ProcControlAPI may deliver callbacks when this function is called.
This function return false on error, and a subsequent call to getLastError
returns details on the error.
}

\begin{apient}
bool stopProc();
\end{apient}

\apidesc{
This function will move all threads in the target process into the stopped
state. This function returns true if at least one thread was stopped as part
of the call, and false otherwise.
It is an error to call this function from a callback.
ProcControlAPI may deliver callbacks when this function is called.
This function return false on error, and a subsequent call to getLastError
returns details on the error.
}

\begin{apient}
bool detach(bool leaveStopped = false);
\end{apient}

\apidesc{
This function will detach ProcControlAPI from the target process.
ProcControlAPI will no longer be able to control or receive events from the
target process. All breakpoints will be removed from the target. This
function returns true on success and false on error. Upon an error a
subsequent call to getLastError returns details on the error.
If the leaveStopped parameter is set to true, and the process is in a stopped
state, then the target process will be left in a stopped state after the
detach.
It is an error to call this function from a callback.
}

\begin{apient}
bool temporaryDetach();
\end{apient}

\apidesc{
This function temporarily detaches from the target process, but leaves the
Process data structure intact. This functionality is commonly called
detach-on-the-fly. The target process will not report new events nor be
controllable or able to be queried by the user. Breakpoints are removed from
the process. The reAttach function will reconnect the process after this
call.
This function returns true on success and false upon error.
It is an error to call this function from a callback.
}

\begin{apient}
bool reAttach();
\end{apient}

\apidesc{
This function reconnects to the target process after a temporaryDetach call.
Any breakpoints will be re-inserted back into the function, and if threads
have been created or destroyed during the time detached new events will be
thrown for them.
This function returns true on success and false upon error.
It is an error to call this function from a callback.
}

\begin{apient}
bool terminate();
\end{apient}

\apidesc{
This function forcefully terminated the target process. Upon a successful call
to this function the target process will end execution. The Process object
will record the target process as having crashed. This function returns true
on success and false on error. Upon an error a subsequent call to
getLastError returns details on the error.
It is an error to call this function from a callback.
}

\begin{apient}
ThreadPool& threads();
const ThreadPool& threads() const;
\end{apient}

\apidesc{
Returns the \code{ThreadPool} of the process.
}

\begin{apient}
const LibraryPool& libraries() const;
LibraryPool& libraries();
\end{apient}

\apidesc{
Returns the \code{LibraryPool} of the process.
}

\begin{apient}
bool addLibrary(std::string libname);
\end{apient}

\apidesc{
This function causes the specified library to be loaded into the process. It
will trigger an event (and thus a user callback) for each library loaded
(including dependencies).
}

\begin{apient}
void* getData() const;
void setData(void* p) const;
\end{apient}

\apidesc{
These functions respectively get and set an opaque data object that can be
associated with this process. The data is not interpreted by ProcControlAPI,
but remains associated with the process.
}

\begin{apient}
unsigned getMemoryPageSize() const;
\end{apient}

\apidesc{
This function returns memory page size for the current OS on which the target
process is running.
}

\begin{apient}
Dyninst::Address mallocMemory(size_t size);
Dyninst::Address mallocMemory(size_t size, Dyninst::Address addr);
\end{apient}

\apidesc{
These functions allocate a region of memory in the target
process' address space of size size. Upon a successful
call these functions will map an area of memory in the target process that is
readable, writeable and executable. The \code{mallocMemory(size\_t)} function will
allocate memory at any available address. The \code{mallocMemory(size\_t, Dyninst::Address)}
function will only allocate memory at the specified address, addr.
It is an error to call this function from a callback.
ProcControlAPI may deliver callbacks when this function is called.
Upon success these functions return the start address of memory that was
allocated and 0 otherwise. Upon an error a subsequent call to getLastError
returns details on the error.
}

\begin{apient}
bool freeMemory(Dyninst::Address addr);
\end{apient}

\apidesc{
This function will free a region of memory that was allocated by the
mallocMemory function. Upon a successful call to this function, the area of
memory starting at addr will be unmapped and no longer accessible to the target
process. It is an error to call this function with an address that was not
returned by mallocMemory.
It is an error to call this function from a callback.
ProcControlAPI may deliver callbacks when this function is called.
Upon success this function returns true, otherwise it returns false. Upon an
error a subsequent call to getLastError returns details on the error.
}

\begin{apient}
bool writeMemory(Dyninst::Address addr, const void* buffer, size_t size) const;
\end{apient}

\apidesc{
This function writes to the target process's memory. The
addr parameter specifies an address in the target process to which
ProcControlAPI should write. The buffer and size parameters specify a region
of controller process memory that will be copied into the target process.
It is an error to call this function on a Process that does not have at least
one Thread in a stopped state.
This function returns true on success and false on error. Upon an error a
subsequent call to getLastError returns details on the error.
}

\begin{apient}
bool readMemory(void* buffer, Dyninst::Address addr, size_t size) const;
\end{apient}

\apidesc{
This function reads from the target process' memory. The
addr and size parameters specify an address in the target process from which
ProcControlAPI should read. The buffer parameter specifies an address in the
controller process where ProcControlAPI should write the copied bytes.
It is an error to call this function on a Process that does not have at least
one Thread in a stopped state.
This function returns true on success and false on error. Upon an error a
subsequent call to getLastError returns details on the error.
}

\begin{apient}
bool getMemoryAccessRights(Dyninst::Address addr, mem_perm& rights);
bool setMemoryAccessRights(Dyninst::Address addr, size_t size, mem_perm rights, mem_perm& oldRights);
\end{apient}

\apidesc{
These functions respectively get and set memory permission at the specified
address, addr. The setMemoryAccessRights function also affects a region of
memory in the target process's address space of size size.
}

\begin{apient}
bool findAllocatedRegionAround(Dyninst::Address addr, MemoryRegion& memRegion);
\end{apient}

\apidesc{
This function finds a region of allocated memory memRegion contains address
addr, and returns true on success, otherwise false.
}

\begin{apient}
bool addBreakpoint(Dyninst::Address addr, Breakpoint::ptr bp) const;
\end{apient}

\apidesc{
This function inserts the Breakpoint specified by bp into the target process at
address addr.
It is an error to call this function on a Process that does not have at least
one Thread in a stopped state.
This function returns true on success and false on error. Upon an error a
subsequent call to getLastError returns details on the error.
}

\begin{apient}
bool rmBreakpoint(Dyninst::Address addr, Breakpoint::ptr bp) const;
\end{apient}

\apidesc{
This function removes the Breakpoint specified by bp at address addr from the
target process.
This function returns true on success and false on error. Upon an error a
subsequent call to getLastError returns details on the error.
}

\begin{apient}
bool postIRPC(IRPC::ptr irpc) const;
\end{apient}

\apidesc{
This function posts the given irpc to the Process. ProcControlAPI selects a
Thread from the Process to run the iRPC and put irpc into that
Thread's queue of posted IRPCs.
Each instance of an IRPC object can be posted at most once. It is an error to
attempt to post a single IRPC object multiple times.
This function returns true on success and false on error. Upon an error a
subsequent call to getLastError returns details on the error.
}

\begin{apient}
bool runIRPCSync(IRPC::ptr irpc);
\end{apient}

\apidesc{
This function posts an irpc, similar to Process::postIRPC; continues the thread
the irpc was posted to; and returns when the irpc has completed running. The
thread will be returned to its original running state when this function
returns.
This function returns true if the irpc was successfully run, and false
otherwise. Note that stopping the thread that is running the irpc while this
function waits for irpc completion causes this function to return an error.
It is an error to call this function from a callback.
}

\begin{apient}
bool runIRPCAsync(IRPC::ptr irpc);
\end{apient}

\apidesc{
This function posts an irpc, similar to Process::postIRPC, and then continues
the thread the irpc was posted to.
This function returns true if the irpc was successfully posted and run, and
false otherwise.
It is an error to call this function from a callback.
}

\begin{apient}
bool getPostedIRPCs(std::vector<IRPC::ptr>& rpcs) const;
\end{apient}

\apidesc{
This function returns all IRPCs posted to this Process in the rpcs vector.
This list does not include any IRPCs currently
running- see Thread::getRunningIRPC() for this
functionality.
This function returns true on success and false on error. Upon an error a
subsequent call to getLastError returns details on the error.
}

\begin{apient}
LibraryTracking* getLibraryTracking();
ThreadTracking* getThreadTracking();
LWPTracking* getLWPTracking();
FollowFork* getFollowFork();
SignalMask* getSignalMask();
const LibraryTracking* getLibraryTracking() const;
const ThreadTracking* getThreadTracking() const;
const LWPTracking* getLWPTracking() const;
const FollowFork* getFollowFork() const;
const SignalMask* getSignalMask() const;
\end{apient}

\apidesc{
These functions return pointers to configuration objects for platform-specific
features associated with this Process object.
These functions return \code{nullptr} if the specified feature is unsupported on the
current platform.
}


    \subsubsection{Memory Permissions}
      The \code{mem\_perm} class represents
general memory page permission for the given memory page in the process.


\begin{apient}
bool Process::mem_perm::read;
bool Process::mem_perm::write;
bool Process::mem_perm::execute;
\end{apient}

\apidesc{
The Process::mem\_perm::read, Process::mem\_perm::write, and
Process::mem\_perm::execute, just as their names imply, respectively represent
read, write, and execution permission of given memory page.
}

\begin{apient}
mem_perm(bool r, bool w, bool x);
\end{apient}

\apidesc{
These constructors provide a convenient way to create the specific memory
permission for the given page.
}

\begin{apient}
bool getR() const;
bool getW() const;
bool getX() const;
\end{apient}

\apidesc{
These functions return true if the given memory page has read, write, and
execution permission, respectively, and false otherwise.
}

\begin{apient}
bool isNone() const;
bool isR()    const;
bool isX()    const;
bool isRW()   const;
bool isRX()   const;
bool isRWX()  const;
\end{apient}

\apidesc{
These functions return true if the permission of given memory page is
NO\_ACCESS, READ\_ONLY, EXECUTE, READ\_WRITE, READ\_EXECUTE, and
READ\_WRITE\_EXECUTE, respectively, and false otherwise.
}

\begin{apient}
Process::mem_perm& setR();
Process::mem_perm& setW();
Process::mem_perm& setX();
\end{apient}

\apidesc{
These functions enable read, write, and execution permission for the given page,
respectively, and return this mem\_perm.
}

\begin{apient}
Process::mem_perm& clrR();
Process::mem_perm& clrW();
Process::mem_perm& clrX();
\end{apient}

\apidesc{
These functions disable read, write, and execution permission for the given
page, respectively, and return this mem\_perm.
}

\begin{apient}
bool operator<(const mem_perm& p) const;
\end{apient}

\apidesc{
This function returns true if this mem\_perm is less than p according to the
notation that read permission encodes to 4, write, 2, and execute, 1, and false
otherwise.
}

\begin{apient}
std::string getPermName() const;
\end{apient}

\apidesc{
Return the memory permission name for this mem\_perm.
}

  \subsection{Class ProcessSet}
    \declaredin{ProcessSet.h}

The ProcessSet class is a set container for multiple
Process objects. It shares many of the same operations as the Process class,
but when an operation is performed on a ProcessSet it is done on every Process
in the ProcessSet. On some systems, a ProcessSet can
achieve better performance when repeating an operation across many target
processes. It satisfies the C++
\href{https://en.cppreference.com/cpp/named_req/ForwardIterator}{LegacyForwardIterator}
concept. It also offers the standard \code{begin/end} members for iteration.

\begin{apient}
  ProcessSet::ptr;
  ProcessSet::const_ptr;
\end{apient}

\apidesc{
  The ptr and const\_ptr types are smart pointers to a ProcessSet object.
}

\begin{apient}
  struct CreateInfo;
  struct AttachInfo;
\end{apient}

\apidesc{
  The CreateInfo and AttachInfo types are used by the
  ProcessSet::createProcessSet and
  ProcessSet::attachProcessSet functions when creating groups of processes.
}

\begin{apient}
  struct write_t;
  struct read_t;
\end{apient}

\apidesc{
  The write\_t and read\_t types are used by ProcessSet::readMemory and
  ProcessSet::writeMemory.
}

\begin{apient}
  static ProcessSet::ptr newProcessSet();
\end{apient}

\apidesc{
  This function creates a new ProcessSet that is empty.
}

\begin{apient}
  static ProcessSet::ptr newProcessSet(Process::const_ptr proc);
\end{apient}

\apidesc{
  This function creates a new ProcessSet containing proc.
}

\begin{apient}
  static ProcessSet::ptr newProcessSet(ProcessSet::const_ptr pset);
\end{apient}

\apidesc{
  This function creates a new ProcessSet that is a copy of pset.
}

\begin{apient}
  static ProcessSet::ptr newProcessSet(const std::set<Process::const_ptr>& procs);
\end{apient}

\apidesc{
  This function creates a new ProcessSet containing every element from procs.
}

\begin{apient}
  static ProcessSet::ptr newProcessSet(AddressSet::const_iterator ab, AddressSet::const_iterator ae);
\end{apient}

\apidesc{
  This function creates a new ProcessSet containing the processes that are found
  within \code{[ab, ae)} of an AddressSet.
}

\begin{apient}
  static ProcessSet::ptr createProcessSet(std::vector<CreateInfo>& cinfo);
\end{apient}

\apidesc{
  This function creates a new ProcessSet by launching new processes. Each
  element in cinfo specifies an executable, arguments, environment and file
  descriptor mappings (with similar semantics to
  Process::createProcess), which are used
  to launch a new process.
  Every successfully created Process will be added to a new ProcessSet that is
  returned by this function.
  In addition, the cinfo vector will be updated so that each
  entry's proc field points to the Process created by that
  entry, and the error\_ret entry will contain an error code for any process
  launch that failed.
}

\begin{apient}
  static ProcessSet::ptr attachProcessSet(std::vector<AttachInfo>& ainfo);
\end{apient}

\apidesc{
  This function creates a new ProcessSet by attaching to existing processes.
  Each element in ainfo specifies a PID and executable (with similar semantics
  to Process::attachProcess), which are
  used to attach to the processes.
  Every successfully attached Process will be added to a new ProcessSet that is
  returned by this function.
  In addition, the ainfo vector will be updated so that each
  entry's proc field points to the Process attached by that
  entry, and the error\_ret entry will contain an error code any process attach
  that failed.
}

\begin{apient}
  ProcessSet::ptr set_union(ProcessSet::ptr pset) const;
\end{apient}

\apidesc{
  This function returns a new ProcessSet whose elements are a set union of this
  ProcessSet and pset.
}

\begin{apient}
  ProcessSet::ptr set_intersection(ProcessSet::ptr pset) const;
\end{apient}

\apidesc{
  This function returns a new ProcessSet whose elements are a set
  intersection of
  this ProcessSet and pset.
}

\begin{apient}
  ProcessSet::ptr set_difference(ProcessSet::ptr pset) const;
\end{apient}

\apidesc{
  This function returns a new ProcessSet whose elements are a set difference of
  this ProcessSet minus pset.
}

\begin{apient}
  iterator find(Process::const_ptr proc);
  const_iterator find(Process::const_ptr proc) const;
\end{apient}

\apidesc{
  These functions search a ProcessSet for the Process pointed to by proc and
  returns an iterator that points to that element. It returns ProcessSet::end()
  if no element is found.
}

\begin{apient}
  iterator find(Dyninst::PID pid);
  const_iterator find(Dyninst::PID pid) const;
\end{apient}

\apidesc{
  These functions search a ProcessSet for the Process pointed to by proc and
  returns an iterator that points to that element. It returns ProcessSet::end()
  if no element is found.
}

\begin{apient}
  bool empty() const;
\end{apient}

\apidesc{
  This function returns true if the ProcessSet has zero elements,
  false otherwise.
}

\begin{apient}
  size_t size() const;
\end{apient}

\apidesc{
  This function returns the number of elements in the ProcessSet.
}

\begin{apient}
  std::pair<iterator, bool> insert(Process::const_ptr proc);
\end{apient}

\apidesc{
  This function inserts proc into the ProcessSet. If proc already exists in the
  ProcessSet, then no change will occur. This function returns an iterator
  pointing to either the new or existing element and a boolean that
  is true if an
  insert happened and false otherwise.
}

\begin{apient}
  void erase(iterator pos);
\end{apient}

\apidesc{
  This function removes the element pointed to by pos from the ProcessSet.
}

\begin{apient}
  size_t erase(Process::const_ptr proc);
\end{apient}

\apidesc{
  This function searches the ProcessSet for proc, then erases that element from
  the ProcessSet. It returns 1 if it erased an element and 0 otherwise.
}

\begin{apient}
  void clear();
\end{apient}

\apidesc{
  This function erases all elements in the ProcessSet.
}

\begin{apient}
  ProcessSet::ptr getErrorSubset() const;
\end{apient}

\apidesc{
  This function returns a new ProcessSet containing every Process from this
  ProcessSet that has a non-zero error code. Error codes are reset upon every
  ProcessSet API call, so this function shows which Processes had an
  error on the
  last ProcessSet operation.
}

\begin{apient}
  void getErrorSubsets(std::map<ProcControlAPI::err_t, ProcessSet::ptr>& err_sets) const;
\end{apient}

\apidesc{
  This function groups every Process from this ProcessSet that has a non-zero
  error code into new ProcessSets, one for each error code. For each error code
  generated by the last ProcessSet API operation an element will be added to
  err\_sets, and every Process that has the same error code will be added to the
  new ProcessSet associated with that error code.
}

\begin{apient}
  bool anyTerminated() const;
  bool allTerminated() const;
\end{apient}

\apidesc{
  These functions respectively return true if any or all processes in this
  ProcessSet are terminated, and false otherwise.
}

\begin{apient}
  bool anyExited() const;
  bool allExited() const;
\end{apient}

\apidesc{
  These functions respectively return true if any or all processes in this
  ProcessSet have exited normally, and false otherwise.
}

\begin{apient}
  bool anyCrashed() const;
  bool allCrashed() const;
\end{apient}

\apidesc{
  These functions respectively return true if any or all processes in this
  ProcessSet have crashed normally, and false otherwise.
}

\begin{apient}
  bool anyDetached() const;
  bool allDetached() const;
\end{apient}

\apidesc{
  These functions respectively return true if any or all processes in this
  ProcessSet have been detached, and false otherwise.
}

\begin{apient}
  bool anyThreadStopped() const;
  bool allThreadsStopped() const;
\end{apient}

\apidesc{
  These functions respectively return true if any or all threads in this
  ProcessSet are stopped, and false otherwise.
}

\begin{apient}
  bool anyThreadRunning() const;
  bool allThreadsRunning() const;
\end{apient}

\apidesc{
  These functions respectively return true if any or all threads in this
  ProcessSet are running, and false otherwise.
}

\begin{apient}
  ProcessSet::ptr getTerminatedSubset() const;
\end{apient}

\apidesc{
  This function returns a new ProcessSet, which is a subset of this ProcessSet,
  and contains every Process that is terminated.
}

\begin{apient}
  ProcessSet::ptr getExitedSubset() const;
\end{apient}

\apidesc{
  This function returns a new ProcessSet, which is a subset of this ProcessSet,
  and contains every Process that has exited normally.
}

\begin{apient}
  ProcessSet::ptr getCrashedSubset() const;
\end{apient}

\apidesc{
  This function returns a new ProcessSet, which is a subset of this ProcessSet,
  and contains every Process that has crashed.
}

\begin{apient}
  ProcessSet::ptr getDetachedSubset() const;
\end{apient}

\apidesc{
  This function returns a new ProcessSet, which is a subset of this ProcessSet,
  and contains every Process that is detached.
}

\begin{apient}
  ProcessSet::ptr getAllThreadRunningSubset() const;
  ProcessSet::ptr getAnyThreadRunningSubset() const;
\end{apient}

\apidesc{
  This function returns a new ProcessSet, which is a subset of this ProcessSet,
  and contains every Process that respectively has any or all threads running.
}

\begin{apient}
  ProcessSet::ptr getAllThreadStoppedSubset() const;
  ProcessSet::ptr getAnyThreadStoppedSubset() const;
\end{apient}

\apidesc{
  This function returns a new ProcessSet, which is a subset of this ProcessSet,
  and contains every Process that respectively has any or all threads stopped.
}

\begin{apient}
  bool continueProcs() const;
\end{apient}

\apidesc{
  This function continues every thread in every process of this ProcessSet,
  similar to Process::continueProc. It
  returns true if every process was successfully continued and false otherwise.
}

\begin{apient}
  bool stopProcs() const;
\end{apient}

\apidesc{
  This function stops every thread in every process of this
  ProcessSet, similar to
  Process::stopProc. It returns true if every process was
  successfully stopped and false otherwise.
}

\begin{apient}
  bool detach(bool leaveStopped = false) const;
\end{apient}

\apidesc{
  This function detaches from every process in this ProcessSet, similar to
  Process::detach. It returns true if every
  process detach was successful and false otherwise.
  If the leaveStopped parameter is set to true, and the processes in this
  ProcessSet are stopped, then the processes will be left in a stopped state
  after the detach.
}

\begin{apient}
  bool terminate() const;
\end{apient}

\apidesc{
  This function terminates every process in this ProcessSet, similar to
  Process::terminate. It returns true if
  every process was successfully terminated and false otherwise.
}

\begin{apient}
  bool temporaryDetach() const;
\end{apient}

\apidesc{
  This function does a temporary detach from every process in this ProcessSet,
  similar to Process::temporaryDetach.
  It returns true if every process was successfully detached and false
  otherwise.
}

\begin{apient}
  bool reAttach() const;
\end{apient}

\apidesc{
  This function reattaches to every process in this ProcessSet, similar to
  Process::reAttach. It returns true if every
  process was successfully reAttached and false otherwise.
}

\begin{apient}
  AddressSet::ptr mallocMemory(size_t sz) const;
\end{apient}

\apidesc{
  This function allocates a block of memory of size sz in each process in this
  ProcessSet. The addresses of the allocations are returned in a new AddressSet
  object.
}

\begin{apient}
  bool mallocMemory(size_t size, AddressSet::ptr location) const;
\end{apient}

\apidesc{
  This function allocates a block of memory of size sz in each process in this
  ProcessSet. The memory will be allocated in each process based on the
  Process/Address pairs in location.
  This function's behavior is undefined if location contains
  processes not included in this ProcessSet.
  This function returns true if every allocation happened without
  error and false
  otherwise.
}

\begin{apient}
  bool freeMemory(AddressSet::ptr addrs) const;
\end{apient}

\apidesc{
  This function frees memory allocated by
  Process::mallocMemory or
  ProcessSet::mallocMemory. The
  AddressSet addrs should contain a list of Process/Address pairs that point to
  the memory that should be freed.
  This function's behavior is undefined if addrs contains
  processes not included in this ProcessSet.
  This function returns true if every free happened without error and false
  otherwise.
}

\begin{apient}
  bool readMemory(AddressSet::ptr addrs,
                  std::multimap<Process::ptr, void*>& result,
                  size_t size) const;
\end{apient}

\apidesc{
  This function reads memory from processes in this ProcessSet. addrs should
  contain the addresses to read memory from. size should be the amount of
  memory read from each process. The results of the memory reads will be
  returned by filling in the result multimap. Each process that is read from
  will have an entry in result along with a malloc allocated buffer containing
  the results of the read.
  It is the ProcControlAPI user's responsibility to free the
  memory buffers returned by this function.
  This function's behavior is undefined if addrs contains
  processes not included in this ProcessSet.
  This function returns true if every read happened without error, and false
  otherwise.
}

\begin{apient}
  bool readMemory(AddressSet::ptr addrs,
                  std::map<void *, ProcessSet::ptr>& result,
                  size_t size,
                  bool use_checksum = true) const;
\end{apient}

\apidesc{
  This function reads memory from processes in this ProcessSet. addrs should
  contain the addresses to read memory from. size should be the amount of
  memory to read from each process. The results of the memory reads will be
  aggregated together into the result map. If any two processes read equivalent
  byte-for-byte data, then those processes are grouped together in a new
  ProcessSet associated with a common malloc allocated buffer containing their
  memory contents.
  It is the ProcControlAPI user's responsibility to free the
  memory buffers returned by this function.
  This function's behavior is undefined if addrs contains
  processes not included in this ProcessSet.
  This function returns true if every read happened without error, and false
  otherwise.
}

\begin{apient}
  bool readMemory(std::multimap<Process::const_ptr, read_t>& addr) const;
\end{apient}

\apidesc{
  This function reads memory from processes in this ProcessSet. The processes to
  read from are specified in the indexes of addr. The remote address, read size
  and local buffer are specified in the read\_t elements of addr.
  This function's behavior is undefined if addr contains
  processes not included in this ProcessSet.
  This function returns true if every read happened without error, and false
  otherwise. If any read results in an error, then the error\_ret field of the
  associated addr element will be set.
}

\begin{apient}
  bool writeMemory(AddressSet::ptr addrs, const void* buffer, size_t sz) const;
\end{apient}

\apidesc{
  This function will write the contents of buffer of size sz into the memory of
  each process at addrs.
  This function's behavior is undefined if addrs contains
  processes not included in this ProcessSet.
  This function returns true if every write happened without error, and false
  otherwise.
}

\begin{apient}
  bool writeMemory(std::multimap<Process::const_ptr, write_t>& addrs) const;
\end{apient}

\apidesc{
  This function writes to the memory of each process in addrs. The processes to
  write to are specified as the indexes of addrs. The local memory buffer,
  buffer size, and target location are specified in the write\_t element of
  addrs.
  This function's behavior is undefined if addrs contains
  processes not included in this ProcessSet.
  This function returns true if every write happened without error, and false
  otherwise. If any write results in an error, then the error\_ret field of the
  associated addr element will be set.
}

\begin{apient}
  bool addBreakpoint(AddressSet::ptr as, Breakpoint::ptr bp) const;
\end{apient}

\apidesc{
  This function inserts the Breakpoint bp into every process and
  address specified
  by as. It is similar to Process::addBreakpoint.
  This function's behavior is undefined if addrs contains
  processes not included in this ProcessSet.
  This function returns true if every breakpoint add happened without error, and
  false otherwise.
}

\begin{apient}
  bool rmBreakpoint(AddressSet::ptr as, Breakpoint::ptr bp) const;
\end{apient}

\apidesc{
  The function removes the Breakpoint bp from each process at the locations
  specified in as. It is similar to Process::rmBreakpoint.
  This function's behavior is undefined if as contains
  processes not included in this ProcessSet.
  This function returns true if every breakpoint remove happened without error,
  and false otherwise.
}

\begin{apient}
  bool postIRPC(const std::multimap<Process::const_ptr, IRPC::ptr>& rpcs) const;
\end{apient}

\apidesc{
  This function posts the IRPC objects specified in rpcs to their associated
  processes in the multimap. It is similar to Process::postIRPC.
  This function's behavior is undefined if rpcs contains
  processes not included in this ProcessSet.
  This function returns true if every post happened without error, and false
  otherwise.
}

\begin{apient}
  bool postIRPC(IRPC::ptr irpc, std::multimap<Process::ptr, IRPC::ptr>* result = NULL) const;
\end{apient}

\apidesc{
  This function makes a copy of irpc for each Process in this
  ProcessSet and posts
  it to that Process. If result is not \code{nullptr}, then the multimap will be filled
  with each newly created IRPC and the Process to which it was posted. It is
  similar to Process::postIRPC.
  This function returns true if every post happened without error, and false
  otherwise.
}

\begin{apient}
  bool postIRPC(IRPC::ptr irpc,
                AddressSet::ptr addrs,
                std::multimap<Process::ptr, IRPC::ptr>* result = NULL) const;
\end{apient}

\apidesc{
  This function makes a copy of irpc and posts it to each Process in
  addrs at the
  given Address. If result is not \code{nullptr}, then the multimap will be filled with
  each newly created IRPC and the Process to which it was posted. It is similar
  to Process::postIRPC.
  This function's behavior is undefined if rpcs contains
  processes not included in this ProcessSet.
  This function returns true if every post happened without error, and false
  otherwise.
}


  \subsection{Class AddressSet}
    \declaredin{ProcessSet.h}

An \code{AddressSet} is a collection of \code{Process} and
\code{Dyninst::Address} tuples. \code{AddressSet} is used by the
ProcessSet\index{ProcessSet: : : :!} and ThreadSet\index{ThreadSet: : : :!}
classes for performing group operations on large numbers of processes. An
\code{AddressSet} might, for example, represent the location of a symbol across
numerous processes, or the location of a buffer in each process where data can
be written or read.  It satisfies the C++
\href{https://en.cppreference.com/cpp/named_req/ForwardIterator}{LegacyForwardIterator}
concept. It also offers the standard \code{begin/end} members for iteration.

\begin{apient}
AddressSet::ptr;
AddressSet::const_ptr;
\end{apient}

\apidesc{
The AddressSet::ptr and AddressSet::const\_ptr respectively represent a pointer
and a const pointer to an AddressSet object. Both pointer types are reference
counted and will cause the underlying AddressSet object to be cleaned when
there are no more references.
}

\begin{apient}
typedef std::pair<Dyninst::Address, Process::ptr> value_type;
\end{apient}

\apidesc{
}

\begin{apient}
static AddressSet::ptr newAddressSet();
\end{apient}

\apidesc{
This function returns a new AddressSet that is empty.
}

\begin{apient}
static AddressSet::ptr newAddressSet(ProcessSet::const_ptr ps, Dyninst::Address addr);
\end{apient}

\apidesc{
This function returns a new AddressSet initialized with the elements from ps
paired with the Address addr.
}

\begin{apient}
static AddressSet::ptr newAddressSet(ProcessSet::const_ptr ps,
                                     std::string library_name,
                                     Dyninst::Offset off = 0);
\end{apient}

\apidesc{
This function returns a new AddressSet initialized with the elements from ps.
The Address element for each process is calculated by looking up the load
address of library\_name in each Process and adding it to off.
}

\begin{apient}
iterator find(Dyninst::Address addr);
const_iterator find(Dyninst::Address addr) const;
\end{apient}

\apidesc{
These functions return an iterator that points to the first element in the
AddressSet with an address of addr. They return end() if no element matches
addr.
}

\begin{apient}
iterator find(Dyninst::Address addr, Process::const_ptr proc);
const_iterator find(Dyninst::Address addr, Process::const_ptr proc) const;
\end{apient}

\apidesc{
These functions return an iterator that points to any element that has a process
and address of proc and addr. It returns end() if no element matches.
}

\begin{apient}
size_t count(Dyninst::Address addr) const;
\end{apient}

\apidesc{
This function returns the number of elements with address addr.
}

\begin{apient}
size_t size() const;
\end{apient}

\apidesc{
This function returns the number of elements in the AddressSet.
}

\begin{apient}
bool empty() const;
\end{apient}

\apidesc{
This function returns true if the AddressSet has zero elements and false
otherwise.
}

\begin{apient}
std::pair<iterator, bool> insert(Dyninst::Address addr, Process::const_ptr proc);
\end{apient}

\apidesc{
This function inserts a new element into the AddressSet with addr and proc as
its values. If another element with those values already exists, then no new
element will be inserted. It returns an iterator that points to the new or
existing element and a boolean value that is true if a new element was inserted
and false otherwise.
}

\begin{apient}
size_t insert(Dyninst::Address addr, ProcessSet::const_ptr ps);
\end{apient}

\apidesc{
For every element in ps, this function inserts it and addr into the AddressSet.
It returns the number of new elements created.
}

\begin{apient}
void erase(iterator pos);
\end{apient}

\apidesc{
This function removes the element pointed to by pos from the AddressSet.
}

\begin{apient}
size_t erase(Process::const_ptr proc);
\end{apient}

\apidesc{
This function removes every element with a process of proc from the AddressSet.
It returns the number of elements removed.
}

\begin{apient}
size_t erase(Dyninst::Address addr, Process::const_ptr proc);
\end{apient}

\apidesc{
This function removes any element that has and address and process of addr and
proc from the AddressSet. It returns the number of elements removed.
}

\begin{apient}
void clear();
\end{apient}

\apidesc{
This function erases all elements from the AddressSet leaving an AddressSet of
size zero.
}

\begin{apient}
iterator lower_bound(Dyninst::Address addr);
const_iterator lower_bound(Dyninst::Address addr) const;
\end{apient}

\apidesc{
This function returns an iterator pointing to the first element in the
AddressSet that has an address greater than or equal to addr.
}

\begin{apient}
iterator upper_bound(Dyninst::Address addr);
const_iterator upper_bound(Dyninst::Address addr) const;
\end{apient}

\apidesc{
This function returns an iterator pointing to the first element in the
AddressSet that has an address greater than addr.
}

\begin{apient}
std::pair<iterator, iterator> equal_range(Address addr) const;
\end{apient}

\apidesc{
This function returns a pair of iterators. The first iterator has the same
value as the return of lower\_bound(addr) and the second iterator has the same
value as the return of upper\_bound(addr).
}

\begin{apient}
AddressSet::ptr set_union(AddressSet::const_ptr aset) const;
\end{apient}

\apidesc{
This function returns a new AddressSet whose elements are the set union of this
AddressSet and aset.
}

\begin{apient}
AddressSet::ptr set_intersection(AddressSet::const_ptr aset) const;
\end{apient}

\apidesc{
This function returns a new AddressSet whose elements are the set intersection
of this AddressSet and aset.
}

\begin{apient}
AddressSet::ptr set_difference(AddressSet::const_ptr aset) const;
\end{apient}

\apidesc{
This function returns a new AddressSet whose elements are the set difference of
this AddressSet minus aset.
}


\section{Threads}
  ProcControlAPI offers several facilities for interacting with threads in a
  platform-independent manner.

  \subsection{Class Thread}
    \declaredin{PCProcess.h}

The \code{Thread} class represents a single thread of execution
in the target process. Any Process has at least one Thread, and
multi-threaded target processes may have more. Each Thread has an associated
integral value known as its LWP, which serves as a handle for communicating
with the OS about the thread (e.g., a PID value on Linux). On some systems,
depending on availability, a Thread may have information from the user space
threading library.

\begin{apient}
Thread::ptr;
Thread::const_ptr;
\end{apient}

\apidesc{
The Thread::ptr and Thread::const\_ptr respectively represent a pointer and a
const pointer to a Thread object. Both pointer types are reference counted
and cause the underlying Thread object to be cleaned when there are no more
references. ProcControlAPI maintains internal references to any Thread it
actively controls, relinquishing those references when the thread exits or is
detached.
}

\begin{apient}
Dyninst::LWP getLWP() const;
\end{apient}

\apidesc{
This function returns an OS handle for this thread. On Linux this returns a
pid\_t for this thread. On FreeBSD, this returns a lwpid\_t.
}

\begin{apient}
Process::ptr getProcess();
Process::const_ptr getProcess() const;
\end{apient}

\apidesc{
These functions return a pointer to the Process object that contains this
thread.
}

\begin{apient}
bool isStopped() const;
\end{apient}

\apidesc{
This function returns true if this thread is in a stopped state and false
otherwise.
}

\begin{apient}
bool isRunning() const;
\end{apient}

\apidesc{
This function returns true if this thread is in a running state and false
otherwise.
}

\begin{apient}
bool isLive() const;
\end{apient}

\apidesc{
This function returns true if this thread is alive, and it returns false if this
thread has been destroyed.
}

\begin{apient}
bool isDetached() const;
\end{apient}

\apidesc{
This function returns true if this thread has been detached via
Process::temporaryDetach and false otherwise.
}

\begin{apient}
bool isInitialThread() const;
\end{apient}

\apidesc{
This function returns true if this thread is the initial thread for the process
and false otherwise.
}

\begin{apient}
bool stopThread();
\end{apient}

\apidesc{
This function moves the thread to into a stopped state. Upon a successful call
to this function the Thread object will be paused and will not resume execution
until the Thread is continued. It is an error to call this function from a
callback. Instead of calling this function, a callback can stop a thread by
returning Process::cbThreadStop or Process::cbProcStop.
ProcControlAPI may deliver callbacks when this function is called.
Upon success this function returns true, otherwise it returns false. Upon an
error a subsequent call to getLastError returns details on the error.
}

\begin{apient}
bool continueThread();
\end{apient}

\apidesc{
This function moves the thread into a running state. It is an error to call
this function from a callback. Instead of calling this function, a callback
can stop a thread by returning Process::cbThreadContinue or
Process::cbProcContinue.
ProcControlAPI may deliver callbacks when this function is called.
Upon success this function returns true, otherwise it returns false. Upon an
error a subsequent call to getLastError returns details on the error.
}

\begin{apient}
bool getRegister(Dyninst::MachRegister reg, Dyninst::MachRegisterVal& val) const;
\end{apient}

\apidesc{
This function gets the value of a single register from this thread. The
register is specified by the reg parameter, and the value of the register is
returned by the val parameter.

It is an error to call this function on a thread that is not in the stopped
state.

Upon success this function returns true, otherwise it returns false. Upon an
error a subsequent call to getLastError returns details on the error.
}

\begin{apient}
bool getAllRegisters(RegisterPool& pool) const;
\end{apient}

\apidesc{
This function reads the values of every register in the thread and returns them
as part of the RegisterPool object pool. Depending on the OS, this call may
be more efficient that calling Thread::getRegister multiple times.
It is an error to call this function on a thread that is not in the stopped
state.
Upon success this function returns true, otherwise it returns false. Upon an
error a subsequent call to getLastError returns details on the error.
}

\begin{apient}
bool setRegister(Dyninst::MachRegister reg, Dyninst::MachRegisterVal val) const;
\end{apient}

\apidesc{
This function writes the value of a single register in this thread. The
register is specified by the reg parameter, and the value that should be
written is specified by the val parameter.

It is an error to call this function on a thread that is not in the stopped
state.
Upon success this function returns true, otherwise it returns false. Upon an
error a subsequent call to getLastError returns details on the error.
}

\begin{apient}
bool setAllRegisters(RegisterPool& pool) const;
\end{apient}

\apidesc{
This function sets the values of every register in this thread to the values
specified in the RegisterPool object pool. Depending on the OS, this call may
be more efficient that calling Thread::setRegister multiple times.
It is an error to call this function on a thread that is not in the stopped
state.
Upon success this function returns true, otherwise it returns false. Upon an
error a subsequent call to getLastError returns details on the error.
}

\begin{apient}
bool haveUserThreadInfo() const;
\end{apient}

\apidesc{
This function returns true if information about this
Thread's underlying user-level thread is available.
}

\begin{apient}
Dyninst::THR_ID getTID() const;
\end{apient}

\apidesc{
This function returns the unique identifier for the user-level thread. This
value is only valid if haveUserThreadInfo returns true.
}

\begin{apient}
Dyninst::Address getStartFunction() const;
\end{apient}

\apidesc{
This function returns the address of the initial function for the user-level
thread. This value is only valid if haveUserThreadInfo returns true.
}

\begin{apient}
Dyninst::Address getStackBase() const;
\end{apient}

\apidesc{
This function returns the address of the bottom of the user-level
thread's stack. This value is only valid if
haveUserThreadInfo returns true.
}

\begin{apient}
unsigned long getStackSize() const;
\end{apient}

\apidesc{
This function returns the size in bytes of the user-level
thread's allocated stack. This value is only valid if
haveUserThreadInfo returns true.
}

\begin{apient}
Dyninst::Address getTLS() const;
\end{apient}

\apidesc{
This function returns the address of the user-level
thread's thread local storage area. This value is only
valid if haveUserThreadInfo returns true.
}

\begin{apient}
bool readThreadLocalMemory(void* buffer,
                           Library::const_ptr lib,
                           Dyninst::Offset tls_symbol_offset,
                           size_t size) const;
\end{apient}

\apidesc{
This function reads from a symbol in thread local storage (TLS) memory. TLS is
memory that is local to a thread and has a lifetime matching the thread. The
tls\_symbol\_offset is the TLS symbol's offset in lib, and
can be found by reading a TLS symbol's value. The lib
parameter can point to a library or the executable. The buffer parameter
specifies an address in the controller process where ProcControlAPI should
write the copied bytes.
It is an error to call this function on a Thread that is not in a stopped state.
It is also an error to call this function on a Thread that has not have
user-level thread information, which can be tested with haveUserThreadInfo.
This function returns true on success and false on error. Upon an error a
subsequent call to getLastError returns details on the error.
}

\begin{apient}
bool writeThreadLocalMemory(Library::const_ptr lib,
                            Dyninst::Offset tls_symbol_offset,
                            const void* buffer,
                            size_t size) const;
\end{apient}

\apidesc{
This function writes to a symbol in thread local storage (TLS) memory. TLS is
memory that is local to a thread and has a lifetime matching the thread. The
tls\_symbol\_offset is the TLS symbol's offset in lib, and
can be found by reading a TLS symbol's value. The lib
parameter can point to a library or the executable. The buffer parameter
specifies an address in the controller process where ProcControlAPI should read
the bytes to be copied.
It is an error to call this function on a Thread that is not in a stopped state.
It is also an error to call this function on a Thread that has not have
user-level thread information, which can be tested with haveUserThreadInfo.
This function returns true on success and false on error. Upon an error a
subsequent call to getLastError returns details on the error.
}

\begin{apient}
bool getThreadLocalAddress(Library::const_ptr lib,
                           Dyninst::Offset tls_symbol_offset,
                           Dyninst::Address& result_addr) const;
\end{apient}

\apidesc{
This function looks up the address of a symbol in thread local storage (TLS)
memory. The tls\_symbol\_offset is the TLS symbol's offset
in lib, and can be found by reading a TLS symbol's value.
The lib parameter can point to a library or the executable. The result\_addr
parameter will be set to the target address for the TLS symbol in this Thread.
It is an error to call this function on a Thread that is not in a stopped state.
It is also an error to call this function on a Thread that has not have
user-level thread information, which can be tested with haveUserThreadInfo.
This function returns true on success and false on error. Upon an error a
subsequent call to getLastError returns details on the error.
}

\begin{apient}
bool postIRPC(IRPC::ptr irpc) const;
\end{apient}

\apidesc{
This function posts the given irpc to the Thread. The IRPC is put irpc into
the Thread's queue of posted IRPCs and will be run when
ready.
Each instance of an IRPC object can be posted at most once. It is an error to
attempt to post a single IRPC object multiple times.
This function returns true on success and false on error. Upon an error a
subsequent call to getLastError returns details on the error.
}

\begin{apient}
bool getPostedIRPCs(std::vector<IRPC::ptr>& rpcs) const;
\end{apient}

\apidesc{
This function returns all IRPCs posted to this Thread in the vector rpcs. This
does not include any running IRPC.
This function returns true on success and false on error. Upon an error a
subsequent call to getLastError returns details on the error.
}

\begin{apient}
IRPC::const_ptr getRunningIRPC() const;
\end{apient}

\apidesc{
This function returns a const pointer to any IRPC that is actively running on
this Thread. If there is no IRPC actively running, then this function returns
IRPC::const\_ptr().
}

\begin{apient}
bool setSingleStepMode(bool mode) const;
\end{apient}

\apidesc{
This function sets whether a Thread is in single-step mode. If called with a
mode of true, then the Thread is put in single-step mode. If called with a
mode of false, then the Thread is taken out of single-step mode.
A Thread in single-step mode will pause execution at each instruction and
trigger an EventSingleStep event. After each EventSingleStep is handled (and
presuming the Thread is still running and in single-step mode) the Thread will
execute one more instruction and trigger another EventSingleStep.
}

\begin{apient}
bool getSingleStepMode() const;
\end{apient}

\apidesc{
This function returns true if the Thread is in single-step mode and false
otherwise.
}

\begin{apient}
void* getData() const;
void setData(void* p) const;
\end{apient}

\apidesc{
These functions respectively get and set an opaque data object that can be
associated with this Thread. The data is not interpreted by ProcControlAPI,
but remains associated with the Thread.
}


  \subsection{Class ThreadPool}
    \declaredin{PCProcess.h}

A \code{ThreadPool} is a collection of threads that make up a
process. Each \code{Process} has one \code{ThreadPool}. A \code{ThreadPool}
is typically used to iterate over or search the set of Threads. It satisfies the C++
\href{https://en.cppreference.com/cpp/named_req/ForwardIterator}{LegacyForwardIterator}
concept. It also offers the standard \code{begin/end} members for iteration.

Note that it is not safe to make assumptions about having consistent contents of
a ThreadPool for a running target process. As the target process runs, threads
may be inserted or removed from the ThreadPool. It is generally safer
to stop a Process before operating on its ThreadPool. When used on a running
process the ThreadPool iterator methods guarantee that they will not return
invalid Thread objects (e.g, nothing that would lead to a segfault), but they
do not guarantee that the Thread objects will refer to live threads or that
they return all Threads.

\begin{apient}
ThreadPool::iterator find(Dyninst::LWP lwp);
ThreadPool::const_iterator find(Dyninst::LWP lwp) const;
\end{apient}

\apidesc{
The functions respectively return an iterator and a const\_iterator that points
to the Thread with a LWP of lwp. If lwp is not found in the thread list, then
this function returns end().
}

\begin{apient}
size_t size() const;
\end{apient}

\apidesc{
This function returns the number of Threads in the Process.
}

\begin{apient}
Process::ptr getProcess();
Process::const_ptr getProcess() const;
\end{apient}

\apidesc{
These functions respectively return a pointer or a const pointer to the Process
that owns this ThreadPool.
}

\begin{apient}
Thread::ptr getInitialThread();
Thread::const_ptr getInitialThread() const;
\end{apient}

\apidesc{
These functions respectively return a pointer or a const pointer to the initial
Thread in a Process. The initial thread is the thread that started execution
of the process (i.e., the thread that called main).
}

  \subsection{Class ThreadSet}
    \declaredin{ProcessSet.h}

A \code{ThreadSet} is a set of Thread pointers. It has similar
operations as Thread, and operations done on a ThreadSet affect every Thread in
that ThreadSet. One some system, such as Blue Gene Q, using a ThreadSet is
more efficient when doing the same operation across a large number of Threads.
It satisfies the C++
\href{https://en.cppreference.com/cpp/named_req/ForwardIterator}{LegacyForwardIterator}
concept. It also offers the standard \code{begin/end} members for iteration.

\begin{apient}
ThreadSet::ptr;
ThreadSet::const_ptr;
\end{apient}

\apidesc{
The ptr and const\_ptr types are smart pointers to a ThreadSet object. When
the last smart pointer to the ThreadSet is cleaned, then the underlying
ThreadSet is cleaned. The const\_ptr type is a const smart pointer.
}

\begin{apient}
ThreadSet::weak_ptr;
ThreadSet::const_weak_ptr;
\end{apient}

\apidesc{

The weak\_ptr and const\_weak\_ptr are weak smart pointers to a ThreadSet
object. Unlike regular smart pointers, weak pointers are not counted as
references when determining whether to clean the ThreadSet object. The
const\_weak\_ptr type is a const weak smart pointer.
}

\begin{apient}
static ThreadSet::ptr newThreadSet();
\end{apient}

\apidesc{
This function creates a new ThreadSet that is empty.
}

\begin{apient}
static ThreadSet::ptr newThreadSet(Thread::ptr thr);
\end{apient}

\apidesc{
This function creates a new ThreadSet that contains thr.
}

\begin{apient}
static ThreadSet::ptr newThreadSet(const ThreadPool& threadp);
\end{apient}

\apidesc{
This function creates a new ThreadSet that contains all of the Threads currently
in threadp.
}

\begin{apient}
static ThreadSet::ptr newThreadSet(const std::set<Thread::const_ptr>& thrds);
\end{apient}

\apidesc{
This function creates a new ThreadSet that contains all of the threads in thrds.
}

\begin{apient}
static ThreadSet::ptr newThreadSet(ProcessSet::ptr pset, bool initial_only = false);
\end{apient}

\apidesc{
This function creates a new ThreadSet that contains every live thread currently
in every process in pset. If initial\_only is true then the new ThreadSet
instead contains only the initial thread of each process in pset.
}

\begin{apient}
ThreadSet::ptr set_union(ThreadSet::ptr tset) const;
\end{apient}

\apidesc{
This function returns a new ThreadSet whose elements are a set union of this
ThreadSet and tset.
}

\begin{apient}
ThreadSet::ptr set_intersection(ThreadSet::ptr tset) const;
\end{apient}

\apidesc{
This function returns a new ThreadSet whose elements are a set intersection of
this ThreadSet and tset.
}

\begin{apient}
ThreadSet::ptr set_difference(ThreadSet::ptr tset) const;
\end{apient}

\apidesc{
This function returns a new ThreadSet whose elements are a set difference of
this ThreadSet minus tset.
}

\begin{apient}
iterator find(Thread::const_ptr thr);
const_iterator find(Thread::const_ptr thr) const;
\end{apient}

\apidesc{
These functions search a ThreadSet for thr and returns an iterator pointing to
that element. It returns ThreadSet::end() if no element is found
}

\begin{apient}
bool empty() const;
\end{apient}

\apidesc{
This function returns true if the ThreadSet has zero elements and false
otherwise.
}

\begin{apient}
size_t size() const;
\end{apient}

\apidesc{
This function returns the number of elements in the ThreadSet.
}

\begin{apient}
std::pair<iterator, bool> insert(Thread::const_ptr thr);
\end{apient}

\apidesc{
This function inserts thr into the ThreadSet. If thr already exists in the
ThreadSet, then no change will occur. This function returns an iterator
pointing to either the new or existing element and a boolean that is true if an
insert happened and false otherwise.
}

\begin{apient}
void erase(iterator pos);
\end{apient}

\apidesc{
This function removes the element pointed to by pos from the ThreadSet.
}

\begin{apient}
size_t erase(Thread::const_ptr thr);
\end{apient}

\apidesc{
This function searches the ThreadSet for thr, then erases that element from the
ThreadSet. It returns 1 if it erased an element and 0 otherwise.
}

\begin{apient}
void clear();
\end{apient}

\apidesc{
This function erases all elements in the ThreadSet.
}

\begin{apient}
ThreadSet::ptr getErrorSubset() const;
\end{apient}

\apidesc{
This function returns a new ThreadSet containing every Thread from this
ThreadSet that has a non-zero error code. Error codes are reset upon every
ThreadSet API call, so this function shows which Threads had an error on the
last ThreadSet operation.
}

\begin{apient}
void getErrorSubsets(std::map<ProcControlAPI::err_t, ThreadSet::ptr>& err) const;
\end{apient}

\apidesc{
This function groups every Thread from this ThreadSet that has a non-zero error
code into new ThreadSets, one for each error code. For each error code generated
by the last ThreadSet API operation an element will be added to err, and every
Thread that has that error code will be added to the new ThreadSet associated
with that error code.
}

\begin{apient}
bool allStopped() const;
bool anyStopped() const;
\end{apient}

\apidesc{
These functions respectively return true if any or all threads in this ThreadSet
are stopped and false otherwise.
}

\begin{apient}
bool allRunning() const;
bool anyRunning() const;
\end{apient}

\apidesc{
These functions respectively return true if any or all threads in this ThreadSet
are running and false otherwise.
}

\begin{apient}
bool allTerminated() const;
bool anyTerminated() const;
\end{apient}

\apidesc{
These functions respectively return true if any or all threads in this ThreadSet
are terminated and false otherwise.
}

\begin{apient}
bool allSingleStepMode() const;
bool anySingleStepMode() const;
\end{apient}

\apidesc{
These functions respectively return true if any or all threads in this ThreadSet
are running in single step mode, and false otherwise.
}

\begin{apient}
bool allHaveUserThreadInfo() const;
bool anyHaveUserThreadInfo() const;
\end{apient}

\apidesc{
These functions respectively return true if any or all threads in this ThreadSet
have user thread information available and false otherwise.
}

\begin{apient}
ThreadSet::ptr getStoppedSubset() const;
\end{apient}

\apidesc{
This function returns a new ThreadSet, which is a subset of this ThreadSet, and
contains every Thread that is stopped.
}

\begin{apient}
ThreadSet::ptr getRunningSubset() const;
\end{apient}

\apidesc{
This function returns a new ThreadSet, which is a subset of this ThreadSet, and
contains every Thread that is running.
}

\begin{apient}
ThreadSet::ptr getTerminatedSubset() const;
\end{apient}

\apidesc{
This function returns a new ThreadSet, which is a subset of this ThreadSet, and
contains every Thread that is terminated.
}

\begin{apient}
ThreadSet::ptr getSingleStepSubset() const;
\end{apient}

\apidesc{
This function returns a new ThreadSet, which is a subset of this ThreadSet, and
contains every Thread that is in single step mode.
}

\begin{apient}
ThreadSet::ptr getHaveUserThreadInfoSubset() const;
\end{apient}

\apidesc{
This function returns a new ThreadSet, which is a subset of this ThreadSet, and
contains every Thread that has user thread information available.
}

\begin{apient}
bool getStartFunctions(AddressSet::ptr result) const;
\end{apient}

\apidesc{
This function fills in the AddressSet pointed to by result with the address of
every start function of each Thread in this ThreadSet. This information is
only available on threads that have user thread information available.
This function return true if it succeeded for every Thread, and false otherwise.
}

\begin{apient}
bool getStackBases(AddressSet::ptr result) const;
\end{apient}

\apidesc{
This function fills in the AddressSet pointed to by result with the address of
every stack base of each Thread in this ThreadSet. This information is only
available on threads that have user thread information available.
This function return true if it succeeded for every Thread, and false otherwise.
}

\begin{apient}
bool getTLSs(AddressSet::ptr result) const;
\end{apient}

\apidesc{
This function fills in the AddressSet pointed to by result with the address of
every thread-local storage region of each Thread in this ThreadSet. This
information is only available on threads that have user thread information
available.
This function return true if it succeeded for every Thread, and false otherwise.
}

\begin{apient}
bool stopThreads() const;
\end{apient}

\apidesc{
This function stops every Thread in this ThreadSet. It is similar to
Thread::stopThread.
This function return true if it succeeded for every Thread, and false
otherwise.
}

\begin{apient}
bool continueThreads() const;
\end{apient}

\apidesc{
This function stops every Thread in this ThreadSet. It is similar to
Thread::continueThread.
This function return true if it succeeded for every Thread, and false
otherwise.
}

\begin{apient}
bool setSingleStepMode(bool v) const;
\end{apient}

\apidesc{
This function puts every Thread in this ThreadSet into single step mode if v is
true. It clears single step mode if v is false. It is similar to
Thread::setSingleStepMode.
This function return true if it succeeded for every Thread, and false
otherwise.
}

\begin{apient}
bool getRegister(Dyninst::MachRegister reg,
                 std::map<Thread::ptr, Dyninst::MachRegisterVal>& res) const;
\end{apient}

\apidesc{
This function gets the value of register reg in every Thread in this ThreadSet.
The collected values are returned in the res map, with each Thread mapped to
the value of reg in that thread. It is similar to
Thread::getRegister.
This function return true if it succeeded for every Thread, and false otherwise.
}

\begin{apient}
bool getRegister(Dyninst::MachRegister reg,
                 std::map<Dyninst::MachRegisterVal, ThreadSet::ptr>& res) const;
\end{apient}

\apidesc{
This function gets the value of register reg in every Thread in this ThreadSet
and then aggregates all identical values together. The res map will contain
an entry for each unique register value, and map that value to a new ThreadSet
that contains every Thread that produced that register value. It is similar to
Thread::getRegister.
This function return true if it succeeded for every Thread, and false otherwise.
}

\begin{apient}
bool setRegister(Dyninst::MachRegister reg,
                 const std::map<Thread::const_ptr, Dyninst::MachRegisterVal>& vals) const;
\end{apient}

\apidesc{
This function sets the value of register reg in each Thread in this ThreadSet.
The value set in each thread is looked up in the vals map. It is similar to
Thread::setRegister.
This function's behavior is undefined if it is passed a
Thread that is not in this ThreadSet.
This function return true if it succeeded for every Thread, and false otherwise.
}

\begin{apient}
bool setRegister(Dyninst::MachRegister reg, Dyninst::MachRegisterVal val) const;
\end{apient}

\apidesc{
This function sets the register reg to val in each Thread in this ThreadSet.
It is similar to Thread::setRegister.
This function return true if it succeeded for every Thread, and false otherwise.
}

\begin{apient}
bool getAllRegisters(std::map<Thread::ptr, RegisterPool>& results) const;
\end{apient}

\apidesc{
This function gets the values of every register in each Thread in this
ThreadSet. The register values are returned as RegisterPools in the results
map, with each Thread mapped to its RegisterPool. It is similar to
Thread::getAllRegisters.
This function return true if it succeeded for every Thread, and false otherwise.
}

\begin{apient}
bool setAllRegisters(const std::map<Thread::const_ptr, RegisterPool>& val) const;
\end{apient}

\apidesc{
This function sets the values of every register in each Thread in this
ThreadSet. The register values are extracted from the val map, with each
Thread specifying its register values via the map. This function is similar
to Thread::setAllRegisters.
This function's behavior is undefined if it is passed a
Thread that is not in this ThreadSet.
This function return true if it succeeded for every Thread, and false otherwise.
}

\begin{apient}
bool postIRPC(const std::multimap<Thread::const_ptr, IRPC::ptr>& rpcs) const;
\end{apient}

\apidesc{
This function posts an IRPC to every Thread in this ThreadSet. The IRPC to
post to each Thread is specified in the rpcs multimap. This function is
similar to Thread::postIRPC.
This function return true if it succeeded for every Thread, and false otherwise.
}

\begin{apient}
bool postIRPC(IRPC::ptr irpc, std::multimap<Thread::ptr, IRPC::ptr>* result = NULL) const;
\end{apient}

\apidesc{
This function posts a copy of irpc to every Thread in this ThreadSet. If
result is not \code{nullptr}, then the new IRPC objects are returned in the result
multimap, with the Thread mapped to the IRPC that was posted there. This
function is similar to Thread::postIRPC.
This function return true if it succeeded for every Thread, and false otherwise.
}


\section{Events}
  ProcControlAPI captures and delivers events from processes and threads.
  An \code{Event} passed to a callback function can be cast to the
  appropriate child class by inspecting the \code{EventType}.

  \subsection{Class Event}
    \declaredin{Event.h}

The Event class represents an instance of an event happening. Each Event
has an EventType that describes the event and pointers to the Process and
Thread that the event occurred on.

The Event class is an abstract class that is never instantiated. Instead,
ProcControlAPI will instantiate children of the Event class, each of which add
information specific to the EventType. For example, an Event representing a
thread creation will have an EventType of ThreadCreate and can be cast into an
EventNewThread for specific information about the new thread.

An event that occurs on a running thread may cause the process, thread, or
neither to stop running until the event has been handled. The specifics of
what is stopped can change between different event types and operating systems.
Each Event describes whether it stopped the associated process or thread with
a SyncType field. The values of this field can be async (the event stopped
neither the process nor thread), sync\_thread (the event stopped its thread),
or sync\_process (the event stopped all threads in the process). A callback
function can choose how to resume or stop a process or thread using its return
value.

\begin{apient}
  enum SyncType { unset, async, sync_thread, sync_process };
\end{apient}

\apidesc{
  The SyncType type is used to describe how a process or thread is stopped by an
  Event. See the above explanation for more details.
}

\begin{apient}
  Thread::const_ptr getThread() const;
\end{apient}

\apidesc{
  This function returns a const pointer to the Thread object that represents the
  thread this event occurred on.
}

\begin{apient}
  Process::const_ptr getProcess() const;
\end{apient}

\apidesc{
  This function returns a const pointer to the Process object that
  represents the
  process this event occurred on.
}

\begin{apient}
  EventType getEventType() const;
\end{apient}

\apidesc{
  This function returns the EventType associated with this Event.
}

\begin{apient}
  SyncType getSyncType() const;
\end{apient}

\apidesc{
  This function returns the SyncType associated with this Event.
}

\begin{apient}
  std::string name() const;
\end{apient}

\apidesc{
  This function returns a human readable name for this Event.
}

\begin{apient}
  EventTerminate::ptr getEventTerminate();
  EventTerminate::const_ptr getEventTerminate() const;
\end{apient}

\apidesc{
}

\begin{apient}
  EventExit::ptr getEventExit();
  EventExit::const_ptr getEventExit() const;
\end{apient}

\apidesc{
}

\begin{apient}
  EventCrash::ptr getEventCrash();
  EventCrash::const_ptr getEventCrash() const;
\end{apient}

\apidesc{
}

\begin{apient}
  EventForceTerminate::ptr getEventForceTerminate();
  EventForceTerminate::const_ptr getEventForceTerminate() const;
\end{apient}

\apidesc{
}

\begin{apient}
  EventExec::ptr getEventExec();
  EventExec::const_ptr getEventExec() const;
\end{apient}

\apidesc{
}

\begin{apient}
  EventStop::ptr getEventStop();
  EventStop::const_ptr getEventStop() const;
\end{apient}

\apidesc{
}

\begin{apient}
  EventBreakpoint::ptr getEventBreakpoint();
  EventBreakpoint::const_ptr getEventBreakpoint() const;
\end{apient}

\apidesc{
}

\begin{apient}
  EventNewThread::ptr getEventNewThread();
  EventNewThread::const_ptr getEventNewThread() const;
\end{apient}

\apidesc{
}

\begin{apient}
  EventNewUserThread::ptr getEventNewUserThread();
  EventNewUserThread::const_ptr getEventNewUserThread() const;
\end{apient}

\apidesc{
}

\begin{apient}
  EventNewLWP::ptr getEventNewLWP();
  EventNewLWP::const_ptr getEventNewLWP() const;
\end{apient}

\apidesc{
}

\begin{apient}
  EventThreadDestroy::ptr getEventThreadDestroy();
  EventThreadDestroy::const_ptr getEventThreadDestroy() const;
\end{apient}

\apidesc{
}

\begin{apient}
  EventUserThreadDestroy::ptr getEventUserThreadDestroy();
  EventUserThreadDestroy::const_ptr getEventUserThreadDestroy() const;
\end{apient}

\apidesc{
}

\begin{apient}
  EventLWPDestroy::ptr getEventLWPDestroy();
  EventLWPDestroy::const_ptr getEventLWPDestroy() const;
\end{apient}

\apidesc{
}

\begin{apient}
  EventFork::ptr getEventFork();
  EventFork::const_ptr getEventFork() const;
\end{apient}

\apidesc{
}

\begin{apient}
  EventSignal::ptr getEventSignal();
  EventSignal::const_ptr getEventSignal() const;
\end{apient}

\apidesc{
}

\begin{apient}
  EventRPC::ptr getEventRPC();
  EventRPC::const_ptr getEventRPC() const;
\end{apient}

\apidesc{
}

\begin{apient}
  EventSingleStep::ptr getEventSingleStep();
  EventSingleStep::const_ptr getEventSingleStep() const;
\end{apient}

\apidesc{
}

\begin{apient}
  EventLibrary::ptr getEventLibrary();
  EventLibrary::const_ptr getEventLibrary() const;
\end{apient}

\apidesc{
  These functions serve as a form of dynamic\_cast. They cast the Event into a
  child type and return the result of that cast. If the Event object is not of
  the appropriate type for the given function, then they return an empty shared
  pointer, \code{ptr()} or \code{const\_ptr()}.
  For example, if an Event was an instance of an EventRPC, then the
  getEventRPC()
  function would cast it to EventRPC and return the resulting value.
}


  \subsection{Class EventType}
    \declaredin{EventType.h}

The \code{EventType} class represents a type of event. Each instance of an
Event happening has one associated EventType, and callback functions can be
registered against EventTypes. Every EventType has a code---an integral value
that identifies the kind of event. Some EventTypes also have a time---Pre, Post
or None---describing when an Event may occur relative to the operation named by
the code. For example, an EventType with a code of Exit and a time of Pre
(written as pre-exit) is associated with an Event that occurs just before a
process exits and its address space is cleaned, while a code of Exit and a time
of Post is associated with an Event that occurs afterwards.

When using EventTypes to register callback functions a special time value of
Any can be used. This signifies that the callback function should trigger for
both Pre and Post time events. ProcControlAPI will never deliver an Event that
has an EventType with a time of Any.

Events and EventTypes are described at a higher level in
Section~\ref{subsubsec:events}, and Table~\ref{tab:event-subclass} lists the
Event subclass used for each EventType.

\begin{apient}
typedef int Code;

enum EventType::Time {
  Pre = 0,
  Post,
  None,
  Any
};
\end{apient}

\apidesc{
The Code and Time types are respectively used to describe the code and time
values of an EventType.
}

\begin{apient}
static const int Error = -1;
static const int Unset = 0;
static const int Exit = 1;
static const int Crash = 2;
static const int Fork = 3;
static const int Exec = 4;
static const int UserThreadCreate = 5;
static const int LWPCreate = 6;
static const int UserThreadDestroy = 7;
static const int LWPDestroy = 8;
static const int Stop = 9;
static const int Signal = 10;
static const int LibraryLoad = 11;
static const int LibraryUnload = 12;
static const int Bootstrap = 13;
static const int Breakpoint = 14;
static const int RPC = 15;
static const int SingleStep = 16;
static const int Library = 17;
static const int ForceTerminate = 18;
static const int ControlAuthority = 20;
static const int AsyncRead = 21;
static const int AsyncWrite = 22;
static const int AsyncReadAllRegs = 23;
static const int AsyncSetAllRegs = 24;
static const int AsyncFileRead = 25;
static const int PreSyscall = 26;
static const int PostSyscall = 27;
\end{apient}

\apidesc{
These constants describe the possible values for an EventType's code, and are
named as EventType::Exit, EventType::Crash and so on. The Error and Unset codes
are for handling error cases and should not be used for callback functions or
be associated with Events.
}

\begin{apient}
static const int Terminate = 400;
static const int ThreadCreate = 401;
static const int ThreadDestroy = 402;
static const int AsyncIO = 403;
\end{apient}

\apidesc{
These codes group several of the codes above. A callback function registered
against one of them is registered against the codes it stands for, and an
EventType carrying the grouping code is never delivered: a callback registered
for Terminate receives Exit, Crash and ForceTerminate events, one registered
for ThreadCreate receives UserThreadCreate and LWPCreate events, and one
registered for ThreadDestroy receives UserThreadDestroy and LWPDestroy events.

AsyncIO is the exception. No handler claims it and it stands for nothing, so
registering a callback for it fails and returns false rather than registering
anything. Register for the individual Async codes instead.
}

\begin{apient}
static const int MaxProcCtrlEvent = 1000;
\end{apient}

\apidesc{
The EventType codes were implemented as an integer (rather than an enum) to
allow users to create custom EventTypes. Any custom EventType should begin at
the MaxProcCtrlEvent value; all smaller values are reserved by ProcControlAPI.
}

\begin{apient}
EventType(Code e);
EventType(Time t, Code e);
EventType();
\end{apient}

\apidesc{
The one parameter constructor constructs an EventType with the given code and a
time of Any. The two parameter constructor constructs an EventType with the
given time and code values. The default constructor constructs an EventType
with an Unset code and a None time value.
}

\begin{apient}
Code code() const;
\end{apient}

\apidesc{
This function returns the code value of the EventType.
}

\begin{apient}
Time time() const;
\end{apient}

\apidesc{
This function returns the time value of the EventType.
}

\begin{apient}
std::string name() const;
\end{apient}

\apidesc{
This function returns a human readable name for this EventType.
}

\begin{apient}
struct eventtype_cmp {
  bool operator()(const EventType& a, const EventType& b) const;
};
\end{apient}

\apidesc{
This type defines a less-than comparison function for EventTypes. While a
comparison of EventTypes does not have a semantic meaning, it is useful for
inserting EventTypes into maps or other STL containers.
}


  \subsection{Class EventBreakpoint}
    \declaredin{Event.h}

An \code{EventBreakpoint} triggers when the target process encounters
a breakpoint
inserted by ProcControlAPI.

Similar to EventStop, EventBreakpoint will explicitly move the thread or process
to a stopped state. Returning cbDefault from a callback function that has
received EventBreakpoint will leave the target process in a stopped state
rather than restore it to the pre-event state.

Associated EventType Code: \code{Breakpoint}

\begin{apient}
  Dyninst::Address getAddress() const;
\end{apient}

\apidesc{
  This function returns the address at which the breakpoint was hit.
}

\begin{apient}
  void getBreakpoints(std::vector<Breakpoint::const_ptr>& b) const;
\end{apient}

\apidesc{
  This function fills b with pointers to the Breakpoints that were hit. Since
  it is possible to insert multiple Breakpoints at the same location, it is
  possible for this function to add more than one Breakpoint to b.
}


  \subsection{Class EventCrash}
    \declaredin{Event.h}

An \code{EventCrash} triggers when a process performs an abnormal exit (e.g.,
crashing on a memory violation). The process that crashed is referenced with
Event's getProcess function.

An EventCrash may be associated with an EventType of pre-crash or post-crash.
Pre-crash means the process has not yet cleaned up its address space, and thus
memory can still be read or written. Post-crash means the process has cleaned up
its address space, memory is no longer accessible.

Associated EventType Code: \code{Crash}

\begin{apient}
int getTermSignal() const;
\end{apient}

\apidesc{
This function returns the signal that caused the process to crash.
}


  \subsection{Class EventExec}
    \declaredin{Event.h}

An \code{EventExec} triggers when a process performs a UNIX-style exec
operation. An \code{EventType} of post-Exec means the process has completed the
exec and setup its new address space. An \code{EventType} of pre-Exec means the
process has not yet torn down its old address space.

Associated \code{EventType} Code: \code{Exec}

\begin{apient}
std::string getExecPath() const;
\end{apient}

\apidesc{
This function returns the file path to the process' new
executable.
}

  \subsection{Class EventExit}
    \declaredin{Event.h}

An \code{EventExit} triggers when a process performs a normal exit (e.g.,
calling the exit function or returning from main). The process that exited is
referenced with Event's getProcess function.

An \code{EventExit} may be associated with an \code{EventType} of pre-exit or
post-exit.
Pre-exit means the process has not yet cleaned up its address space, and thus
memory can still be read or written. Post-exit means the process has cleaned up
its address space, memory is no longer accessible.

Associated \code{EventType} Code: \code{Exit}

\begin{apient}
int getExitCode() const;
\end{apient}

\apidesc{
This function returns the process' exit code.
}

  \subsection{Class EventForceTerminate}
    \declaredin{Event.h}

An \code{EventForceTerminate} triggers when a process is forcefully terminated
via \code{Process::terminate}. When the callback is delivered for this event,
the address space of the corresponding process will no longer be available.

Associated EventType Code: \code{ForceTerminate}

\begin{apient}
int getTermSignal() const;
\end{apient}

\apidesc{
This function returns the signal that was used to terminate the process.
}


  \subsection{Class EventFork}
    \declaredin{Event.h}

An \code{EventFork} triggers when a process performs a UNIX-style fork
operation. The process that performed the initial fork is returned via Event's
getProcess member function, while the newly created process can be found via
EventFork's getChildProcess member function.

Associated \code{EventType} Code: \code{Fork}

\begin{apient}
Process::const_ptr getChildProcess() const;
\end{apient}

\apidesc{
This function returns a pointer to the Process object that represents the newly
created child process.
}

  \subsection{Class EventLWPDestroy}
    \declaredin{Event.h}

An \code{LWPThreadDestroy} triggers when a thread exits. Event's getThread
member function returns the thread that exited.

Associated EventType Code: \code{LWPDestroy}


  \subsection{Class EventLibrary}
    \declaredin{Event.h}

An \code{EventLibrary} triggers when the process either loads or unloads a
shared library. ProcControlAPI will not trigger an \code{EventLibrary} for
library unloads associated with a Process being terminated, and it will not
trigger \code{EventLibrary} for library loads that happened before a
ProcControlAPI attach operation.

It is possible for multiple libraries to be loaded or unloaded at the same time.
In this case, an \code{EventLibrary} will contain multiple libraries in its load
and unload sets.

Associated \code{EventType} Code: \code{Library}

\begin{apient}
const std::set<Library::ptr>& libsAdded() const;
\end{apient}

\apidesc{
This function returns the set of Library objects that were loaded into the
target process' address space. The set will be empty if
no libraries were loaded.
}

\begin{apient}
const std::set<Library::ptr>& libsRemoved() const;
\end{apient}

\apidesc{
This function returns the set of libraries that were unloaded from the target
process' address space. The set will be empty if no
libraries were unloaded.
}

  \subsection{Class EventNewLWP}
    \declaredin{Event.h}

An \code{EventNewLWP} triggers when a process spawns a new LWP. The Event
class' getThread function returns the original Thread that performed the spawn
operation, while EventNewThread's getNewThread returns the newly created
Thread.

A callback function that receives an EventNewThread can use the two field form
of \code{Process::cb\_ret\_t} to control the parent and child thread.

Associated EventType Code: \code{LWPCreate}

\begin{apient}
Thread::const_ptr getNewThread() const;
\end{apient}

\apidesc{
This function returns a const pointer to the Thread object that represents the
newly spawned thread.
}


  \subsection{Class EventNewThread}
    \declaredin{Event.h}

An \code{EventNewThread} triggers when a process spawns a new thread. The Event
class' getThread function returns the original Thread that performed the spawn
operation, while \code{EventNewThread}'s getNewThread returns the newly created
Thread.

This event is never instantiated by ProcControlAPI and simply serves as a
place-holder type for a user to deal with thread creation without having to deal
with the specifics of LWP and user thread creation.

A callback function that receives an \code{EventNewThread} can use the two field
form of \code{Process::cb\_ret\_t}  to control the parent and child thread.

Associated EventType Codes: \code{ThreadCreate}, \code{UserThreadCreate}, and
\code{LWPCreate}

\begin{apient}
Thread::const_ptr getNewThread() const;
\end{apient}

\apidesc{
This function returns a const pointer to the Thread object that represents the
newly spawned thread.
}


  \subsection{Class EventNewUserThread}
    \declaredin{Event.h}

An \code{EventNewUserThread} triggers when a process spawns a new user-level
thread.
The Event class' getThread function returns the original Thread that performed
the spawn operation. This thread may have already been created if the platform
supports the EventNewLWP event. If not, the getNewThread function returns the
newly created Thread.

A callback function that receives an EventNewThread can use the two field form
of \code{Process::cb\_ret\_t}  to control the parent and child thread.

Associated EventType Code: \code{UserThreadCreate}

\begin{apient}
Thread::const_ptr getNewThread() const;
\end{apient}

\apidesc{
This function returns a const pointer to the Thread object that represents the
newly spawned thread or the corresponding thread, if the thread has already
been created.
}


  \subsection{Class EventPostSyscall}
    \declaredin{Event.h}

An \code{EventPreSyscall} is triggered when a thread enters a system call,
provided that the thread is in system call tracing mode.

Associated EventType Code: \code{Syscall}

\begin{apient}
Dyninst::Address getAddress() const;
\end{apient}

\apidesc{
This function returns the address where the system call occurred.
}

\begin{apient}
MachSyscall getSyscall() const;
\end{apient}

\apidesc{
This function returns information about the system call.
}

  \subsection{Class EventPreSyscall}
    \declaredin{Event.h}

An EventPreSyscall is triggered when a thread enters a system call, provided
that the thread is in system call tracing mode.

Associated EventType Code: \code{Syscall}

\begin{apient}
Dyninst::Address getAddress() const;
\end{apient}

\apidesc{
This function returns the address where the system call occurred.
}

\begin{apient}
MachSyscall getSyscall() const;
\end{apient}

\apidesc{
This function returns information about the system call.
}

  \subsection{Class EventRPC}
    \declaredin{Event.h}

An \code{EventRPC} triggers when a process or thread completes a ProcControlAPI
iRPC. When a callback function receives an EventRPC, the memory and registers
that were used by the iRPC can still be found in the address space and thread
that the iRPC ran on. Once the callback function completes, the registers and
memory are restored to their original state.

Associated EventType Code: \code{RPC}

\begin{apient}
IRPC::const_ptr getIRPC() const;
\end{apient}

\apidesc{
This function returns a const pointer to the IRPC object that completed.
}


  \subsection{Class EventSignal}
    \declaredin{Event.h}

An \code{EventSignal} triggers when a process receives a UNIX style signal.

Associated EventType Code: \code{Signal}

\begin{apient}
int getSignal() const;
\end{apient}

\apidesc{
This function returns the signal number that triggered the EventSignal.
}


  \subsection{Class EventSingleStep}
    \declaredin{Event.h}

An \code{EventSingleStep} triggers when a thread, which was put in single-step
mode by \code{Thread::setSingleStep}, completes a single step operation.
The Thread will remain in single-step mode after completion of this event
(presuming it has not be explicitly disabled by Thread::setSingleStep).

Associated EventType Code: \code{SingleStep}


  \subsection{Class EventStop}
    \declaredin{Event.h}

An \code{EventStop} is triggered when a process is stopped by a
non-ProcControlAPI source. On UNIX based systems, this is triggered by receipt
of a SIGSTOP signal.

Unlike most other events, an EventStop will explicitly move the associated
thread or process (see the Event's SyncType to tell which) to a stopped state.
Returning cbDefault from a callback function that has received EventStop will
leave the target process in a stopped state rather than restore it to the
pre-event state.

Associated EventType Code: \code{Stop}


  \subsection{Class EventTerminate}
    \declaredin{Event.h}

The \code{EventTerminate} class is a parent class for EventExit and EventCrash.
It is never instantiated by ProcControlAPI and simply serves as a place-holder
type for a user to deal with process termination without dealing with the
specifics of whether a process exited properly or crashed.

Associated EventType Codes: \code{Exit}, \code{Crash}, and \code{ForceTerminate}


  \subsection{Class EventThreadDestroy}
    \declaredin{Event.h}

An \code{EventThreadDestroy} triggers when a thread exits. Event's getThread
member function returns the thread that exited.

This event is never instantiated by ProcControlAPI and simply serves as a
place-holder type for a user to deal with thread destruction without having to
deal with the specifics of LWP and user thread destruction.

Associated EventType Codes: \code{ThreadDestroy}, \code{UserThreadDestroy}, and
\code{LWPDestroy}


  \subsection{Class EventUserThreadDestroy}
    \declaredin{Event.h}

An \code{EventUserThreadDestroy} triggers when a thread exits. Event's
getThread member function returns the thread that exited.

If the platform also supports EventLWPDestroy events, this event will precede an
EventLWPDestroy event.

Associated EventType Code: \code{UserThreadDestroy}


  \subsection{Class EventNotify}
    \declaredin{PCProcess.h}

The \code{EventNotify} class is used to notify the ProcControlAPI user when
ProcControlAPI is ready to deliver a callback function. EventNotify is a
singleton class, which means only one instance of it is ever instantiated; that
instance is reached through the evNotify function. Notification is described at
a higher level in Section~\ref{subsubsec:callback-delivery}.

\begin{apient}
typedef void (*notify_cb_t)();
\end{apient}

\apidesc{
This function signature is used for the light-weight notification callback.
}

\begin{apient}
EventNotify* evNotify();
\end{apient}

\apidesc{
This function returns the singleton instance of the EventNotify class. It is a
global function rather than a member of EventNotify.
}

\begin{apient}
int getFD();
\end{apient}

\apidesc{
This function returns a file descriptor. ProcControlAPI will write a byte that
will be available for reading on this file descriptor when a callback function
is ready to be invoked. Upon seeing that a byte has been written to this file
descriptor (likely via select or poll) the user should call the
Process::handleEvents function. The user should never actually read the byte
from this file descriptor; ProcControlAPI will handle clearing the byte after
the callback function is invoked.

This function returns -1 on error. Upon an error a subsequent call to
getLastError returns details on the error.
}

\begin{apient}
void registerCB(notify_cb_t cb);
\end{apient}

\apidesc{
This function registers a light-weight callback function that will be invoked
when ProcControlAPI wishes to notify the user that a callback function is ready
to be invoked. This light-weight callback may be called by a ProcControlAPI
internal thread or from a signal handler; the user is encouraged to keep its
implementation appropriately safe for these circumstances.
}

\begin{apient}
void removeCB(notify_cb_t cb);
\end{apient}

\apidesc{
This function removes a light-weight callback cb that was previously
registered with the EventNotify::registerCB function. ProcControlAPI will no
longer invoke cb after this function completes.
}


\section{Inferior Remote Procedure Calls}

  \subsection{Class IRPC}
    \declaredin{PCProcess.h}

\code{IRPC} represents an Inferior Remote Procedure Call that can be run in a
target process. See \code{Process::postIRPC} and \code{Thread::postIRPC} for
information about posting an IRPC.

\begin{apient}
IRPC::ptr;
IRPC::const_ptr;
\end{apient}

\apidesc{
The IRPC::ptr and IRPC::const\_ptr respectively represent a pointer and a const
pointer to an IRPC object. Both pointer types are reference counted and will
cause the underlying IRPC object to be cleaned when there are no more
references. ProcControlAPI will maintain internal references to any IRPC
currently posted or executing.
}

\begin{apient}
static IRPC::ptr createIRPC(void* binary_blob,
                            unsigned int size,
                            bool non_blocking = false);
static IRPC::ptr createIRPC(void* binary_blob,
                            unsigned int size,
                            Dyninst::Address addr,
                            bool non_blocking = false);
\end{apient}

\apidesc{
The createIRPC static function creates and returns a new IRPC object. The
binary\_blob and size parameters specify a buffer of machine code bytes that
this IRPC should execute when invoked. ProcControlAPI will maintain its own
copy of the binary\_blob buffer, the ProcControlAPI user can free the buffer
once this function completes.
If the non\_blocking parameter is false, its default, then calls to
Process::handleEvents will block until this IRPC has completed. If it is
true then the IRPC runs asynchronously and Process::handleEvents does not
wait for it.
If the addr parameter is given, then ProcControlAPI will write and run the
binary code at addr. Otherwise ProcControlAPI will select a location at
which to run the IRPC.
}

\begin{apient}
Dyninst::Address getAddress() const;
\end{apient}

\apidesc{
The getAddress function returns the address at which the IRPC will be run. If
the IRPC was not given an address at construction and has not yet started
running, then this function may return 0.
}

\begin{apient}
void* getBinaryCode() const;
\end{apient}

\apidesc{
The getBinaryCode function returns ProcControlAPI's copy of the binary code
blob this IRPC was created with. That copy is made when the IRPC is created
and is freed with it, so the returned buffer belongs to ProcControlAPI and
should not be freed by the user.
}

\begin{apient}
unsigned int getBinaryCodeSize() const;
\end{apient}

\apidesc{
The getBinaryCodeSize function returns the size of the binary code blob buffer.
}

\begin{apient}
unsigned long getID() const;
\end{apient}

\apidesc{
The getID function returns an integer identifier that uniquely identifies this
IRPC.
}

\begin{apient}
void setStartOffset(unsigned long off);
\end{apient}

\apidesc{
By default an IRPC will start executing its code blob at the beginning of the
blob. This function can be used to tell ProcControlAPI to start execution of
the code blob at some byte offset, off, into the blob.
This function should be called before the IRPC is posted.
}

\begin{apient}
unsigned long getStartOffset() const;
\end{apient}

\apidesc{
If a start offset has been set for this IRPC, then getStartOffset returns it.
Otherwise this function returns 0.
}


\section{Breakpoints}
  A breakpoint is a point in the code region of a target process that,
  when executed, stops the process execution and notifies ProcControlAPI.
  Upon being continued the process will resume execution at the point. A
  Breakpoint object is a handle that can represents one or more breakpoints in
  one or more processes. Upon receiving notification that a breakpoint has
  executed, ProcControlAPI will deliver a callback with an EventBreakpoint.

  Some Breakpoint objects can be created as control-transfer
  breakpoints. When a process is
  continued after executing a control-transfer the process will resume at an
  alternate location, rather than at the breakpoint's
  installation point.

  A single Breakpoint can be inserted into multiple locations within a target
  process. This can be useful when a user has wants to perform a single action
  at multiple locations in a target process. For example, if a user wants to
  insert a breakpoint at the entry to function foo, and foo has multiple
  instantiations in a process, then a single Breakpoint can be inserted at each
  instance of foo.

  A single Breakpoint object can be inserted into multiple target processes at the
  same time. When a process does an operation that copies an address space,
  such as fork on UNIX, the child process will receive all Breakpoint objects
  that were installed in the parent process.

  Multiple Breakpoint objects can be inserted into the same location with-in the
  same process. When this location is executed in the target process a single
  callback will be delivered, and the EventBreakpoint object will contain a
  reference to each Breakpoint inserted at the location. At most one
  control-transfer breakpoint can be inserted at any one point in a process.

  Due to the many-to-many nature of Breakpoints and Processes, a single
  installation of a Breakpoint can be identified by a Breakpoint, Process,
  Address triple. The functions for inserting and removing breakpoints
  (Process::addBreakpoint and Process::rmBreakpoint) need all three pieces of
  information.

  A Breakpoint can be a hardware breakpoint or
  a software breakpoint. A hardware
  breakpoint is typically implemented by setting special debug register in the
  process and can trigger on code execution, data reads or data write. A
  software breakpoint is typically implemented by writing a special instruction
  into a code sequence and can only be triggered by code execution. There are
  typically a limited number of hardware breakpoints available at the same time.

  \subsection{Class Breakpoint}
    \declaredin{PCProcess.h}

\begin{apient}
Breakpoint::ptr;
Breakpoint::const_ptr;
\end{apient}

\apidesc{
The Breakpoint::ptr and Breakpoint::const\_ptr types are respectively a pointer
and a const pointer to a Breakpoint object. These pointers are shared
pointers, and the underlying Breakpoint object will be automatically clean when
there are no more references to it. ProcControlAPI will automatically
maintain at least one reference to any Breakpoint that is installed in a target
process.
}

\begin{apient}
static const int BP_X = 1;
static const int BP_W = 2;
static const int BP_R = 4;
\end{apient}

\apidesc{
These constant values are used to set execute, write and read bits on hardware
breakpoints.
}

\begin{apient}
static Breakpoint::ptr newBreakpoint();
\end{apient}

\apidesc{
This function creates a new software Breakpoint object and returns it. The
Breakpoint is not inserted into a Process until it is passed to
Process::addBreakpoint().
}

\begin{apient}
static Breakpoint::ptr newTransferBreakpoint(Dyninst::Address ctrl_to);
\end{apient}

\apidesc{
This function creates a new control transfer software breakpoint. Upon
resumption after executing this Breakpoint, control will resume at the address
specified by the ctrl\_to parameter.
}

\begin{apient}
static Breakpoint::ptr newHardwareBreakpoint(unsigned int mode, unsigned int size);
\end{apient}

\apidesc{
This function creates a new hardware breakpoint. The mode parameter is a
bitfield that contains an OR combination the values BP\_X, BP\_W and BP\_R.
These control whether the breakpoint will trigger when its target address is
executed, written or read.
The size parameter specifies a range of memory that this breakpoint monitors.
If memory is accessed between the target address and target address + size,
then the breakpoint will trigger.
}

\begin{apient}
bool isCtrlTransfer() const;
\end{apient}

\apidesc{
This function returns true if the Breakpoint is a control transfer breakpoint,
and false if it is a regular Breakpoint.
}

\begin{apient}
Dyninst::Address getToAddress() const;
\end{apient}

\apidesc{
If this Breakpoint is a control transfer breakpoint, then this function returns
the address to which it transfers control. If this Breakpoint is not a
control transfer breakpoint, then this function returns 0.
}

\begin{apient}
void setData(void* data) const;
\end{apient}

\apidesc{
This function sets the value of an opaque handle that is associated with each
Breakpoint. The opaque handle can be any value, and it can be retrieved with
the getData function.
}

\begin{apient}
void* getData() const;
\end{apient}

\apidesc{
This function returns the value of the opaque handled that is associated with
this Breakpoint.
}

\begin{apient}
void setSuppressCallbacks(bool val);
\end{apient}

\apidesc{
This function can be used to suppress callbacks stemming from a specific
breakpoint when called with val set to true value. All other effects from
this breakpoint will still occur, but it will not generate a callback. By
default callbacks occur from every breakpoint.
}

\begin{apient}
bool suppressCallbacks() const;
\end{apient}

\apidesc{
This function returns true if callbacks have been suppressed for this breakpoint
from Breakpoint::setSuppressCallbacks and false otherwise.
}


\section{Registers}

  \subsection{Class MachRegister}
    \declaredin{MachRegister.h}

The MachRegister interface is used for accessing registers in
ProcControlAPI. The entire definition of MachRegister contains more register
names than are listed here; this appendix only lists the registers that can be
accessed through ProcControlAPI.

An instance of MachRegister is defined for each register ProcControlAPI can
name. These instances live inside a namespace that represents the
register's architecture. For example, we can name a
register from an AMD64 machine with Dyninst::x86\_64::rax or a register from a
Power machine with Dyninst::ppc32::r1.


\begin{apient}
typedef unsigned long MachRegisterVal;
\end{apient}

\apidesc{
The MachRegisterVal type is used to represent the contents of a register. If a
register's contents are smaller than MachRegisterVal, then
it will be up cast into a MachRegisterVal.
}

\begin{apient}
typedef enum { Arch_none, Arch_x86, Arch_x86_64, Arch_ppc32, Arch_ppc64 } Architecture;
\end{apient}

\apidesc{
The Architecture enum represents a system's architecture.
}

\begin{apient}
unsigned int getArchAddressWidth(Architecture arch);
\end{apient}

\apidesc{
Returns the size of a pointer, in bytes, on the given architecture, arch.
}

\begin{apient}
static MachRegister getPC(Dyninst::Architecture arch);
\end{apient}

\apidesc{
Returns the program counter (PC) for the architecture \code{arch}.
}

\begin{apient}
static MachRegister getFramePointer(Dyninst::Architecture arch);
\end{apient}

\apidesc{
Returns the frame pointer for the architecture \code{arch}.
}

\begin{apient}
static MachRegister getStackPointer(Dyninst::Architecture arch);
\end{apient}

\apidesc{
Returns the stack pointer for the architecture \code{arch}.
}

\begin{apient}
MachRegister getBaseRegister() const;
\end{apient}

\apidesc{
This function returns the largest register that may alias with the given
register. For example, getBaseRegister on x86\_64::ah returns x86\_64::rax.
}

\begin{apient}
Architecture getArchitecture() const;
\end{apient}

\apidesc{
This function returns the Architecture for this register.
}

\begin{apient}
bool isValid() const;
\end{apient}

\apidesc{
This function returns true if this register is valid. Some API functions may
return invalid registers upon error.
}

\begin{apient}
const char* name() const;
\end{apient}

\apidesc{
This function returns a human readable name that identifies this register.
}

\begin{apient}
unsigned int size() const;
\end{apient}

\apidesc{
This function returns the size of the register in bytes.
}

\begin{apient}
int32_t val() const;
\end{apient}

\apidesc{
This function returns a unique integer that represents this register. This can
be useful for writing switch statements that take a MachRegister. The unique
integers for a MachRegister can be found by preceding the register object name
with an ``i''. For example, a switch statement for MachRegister, reg, could be written as:
}


\begin{lstlisting}
switch (reg.val())  {
  case x86_64::irax:
  case x86_64::irbx:
  case x86_64::ircx:
    ...
}
\end{lstlisting}

\begin{apient}
bool isPC() const;
\end{apient}

\apidesc{
Returns \code{true} if the register is the program counter (PC) for any architecture.
}

\begin{apient}
bool isFramePointer() const;
\end{apient}

\apidesc{
Returns \code{true} if the register is the frame pointer for any architecture.
}

\begin{apient}
bool isStackPointer() const;
\end{apient}

\apidesc{
Returns \code{true} if the register is the stack pointer for any architecture.
}


  \subsection{Class RegisterPool}
    \declaredin{PCProcess.h}

A \code{RegisterPool} represents a collection of registers.
It can be used to get or set all registers in a Thread at once.
It satisfies the C++
\href{https://en.cppreference.com/cpp/named_req/ForwardIterator}{LegacyForwardIterator}
concept. It also offers the standard \code{begin/end} members for iteration.

\begin{apient}
RegisterPool::iterator find(Dyninst::MachRegister r);
RegisterPool::const_iterator find(Dyninst::MachRegister r) const;
\end{apient}

\apidesc{
This function returns an iterator that points to the element in the register
pool that equals register r. If r is not found, then this function returns
end().
}

\begin{apient}
size_t size() const;
\end{apient}

\apidesc{
This function returns the number of elements in the register set.
}

\begin{apient}
Dyninst::MachRegisterVal& operator[](Dyninst::MachRegister r);
const Dyninst::MachRegisterVal& operator[](Dyninst::MachRegister r) const;
\end{apient}

\apidesc{
This function returns a reference to the value associated with the register r in
this register pool. It can be used to efficiently get and set the values of
registers in this pool, or to fill the pool with new
MachRegisters.
}

\section{System Calls}

  System calls are currently only supported on Linux.

  \subsection{Class MachSyscall}
    \declaredin{MachSyscall.h}

A \code{MachSyscall} represents a system call in a platform-independent manner.

\begin{apient}
SyscallIDPlatform num() const;
\end{apient}

\apidesc{
Returns the platform-specific syscall number
}

\begin{apient}
SyscallName name() const;
\end{apient}

\apidesc{
Returns the name of the system call (e.g. ``getpid'')
}

\begin{apient}
bool operator==(const MachSyscall&) const;
\end{apient}

\apidesc{
Equality operator. Two system calls are equal if they are for the same platform
and have the same syscall number.
}

\begin{apient}
static MachSyscall makeFromPlatform(Platform, SyscallIDIndependent);
\end{apient}

\apidesc{
Constructs a MachSyscall from a Platform and a platform-independent ID (e.g.
dyn\_getpid).
}

\section{Libraries}
  ProcControlAPI offers several facilities for interacting with libraries
  both in a created address spaces as well as loading new libraries into one.

  \subsection{Class Library}
    \declaredin{PCProcess.h}

A Library represents a single shared library (frequently
referred to as a DLL or DSO, depending on the OS) that has been loaded into the
target process. In addition, a Library will be used to represent the
process' executable. Processes with statically linked executables
will only contain the single Library that
represents the executable.

Each Library contains a load address and a file
name. The load address is the address at which the OS loaded the library, and
the file name is the path to the library's file. Note
that on some operating systems (e.g., Linux and FreeBSD) the load
address does not necessarily represent the beginning of the library in memory;
instead it is a value that can be added to a library's
symbol offsets to compute the dynamic address of a symbol.

Libraries may be loaded and unloaded by the process during execution. A
library load or unload can trigger a callback with an EventLibrary parameter.
The current list of libraries loaded into a process can be accessed via the
\code{LibraryPool} of a \code{Process}.

\begin{apient}
  Library::ptr;
  Library::const_ptr;
\end{apient}

\apidesc{
  \code{typedef}s for a pointer and a const pointer to a library.
}

\begin{apient}
  std::string getName() const;
\end{apient}

\apidesc{
  Returns the file name for this Library.
}

\begin{apient}
  std::string getAbsoluteName() const;
\end{apient}

\apidesc{
  Returns a file name for this Library that does not contain symlinks or a
  relative path.
}

\begin{apient}
  Dyninst::Address getLoadAddress() const;
\end{apient}

\apidesc{
  Returns the load address for this Library.
}

\begin{apient}
  Dyninst::Address getDataLoadAddress() const;
\end{apient}

\apidesc{
  The AIX operating system can have two load addresses for a library:
  one for the
  code region and one for the data region. On AIX Library::getLoadAddress
  returns the load address of the code region and Library::getDataLoadAddress
  returns the load address of the data region. On non-AIX systems this function
  returns 0.
}

\begin{apient}
  Dyninst::Address getDynamicAddress() const;
\end{apient}

\apidesc{
  On ELF based systems (FreeBSD, Linux) this function will return the
  address of the Library's dynamic section. On other
  systems this function returns 0.
}

\begin{apient}
  bool isSharedLib() const;
\end{apient}

\apidesc{
  Returns true of this Library object is refers to a shared library
  and false if it refers to an executable.
}

\begin{apient}
  void* getData() const;
\end{apient}

\apidesc{
  Returns an opaque data object that user code can associate with
  this library. Use setData to set this opaque value.
}

\begin{apient}
  void setData(void* p) const;
\end{apient}

\apidesc{
  This function sets associates an opaque data object with the Library.
  ProcControlAPI does not try to interpret this value, but will return it with
  the getData function.
}


  \subsection{Class LibraryPool}
    \declaredin{PCProcess.h}

A \code{LibraryPool} is a container representing the executable
and set shared libraries (e.g., .dll and .so libraries) loaded into the target
process' address space. A statically linked target
process will only have a single executable, while a dynamically linked target
process will have an executable and zero or more shared libraries. It satisfies
the C++
\href{https://en.cppreference.com/cpp/named_req/ForwardIterator}{LegacyForwardIterator}
concept. It also offers the standard \code{begin/end} members for iteration.

\begin{apient}
  size_t size() const;
\end{apient}

\apidesc{
  This function returns the number of elements in the library set.
}

\begin{apient}
  Library::ptr getExecutable();
  Library::const_ptr getExecutable() const;
\end{apient}

\apidesc{
  These functions respectively return a pointer or a const pointer to
  the Library object that represents the target process' executable.
}

\begin{apient}
  Library::ptr getLibraryByName(std::string name);
  Library::const_ptr getLibraryByName(std::string name) const;
\end{apient}

\apidesc{
  These functions respectively return a pointer or a const pointer to
  the Library object that with a file name equal to name. If no library
  is found, returns an iterator that compares equal with \code{end()}
}


\section{Platform-specific Features}
  The classes described in this section are all used to configure
  platform-specific features for Process objects. The three tracking classes
  (LibraryTracking, ThreadTracking, LWPTracking) all contain member functions to
  set either interrupt-driven or polling-driven handling for different events
  associated with Process objects. When interrupt-driven handling is enabled,
  the associated process may be modified to accommodate timely handling (e.g.,
  inserting breakpoints). When polling-driven handling is enabled, the
  associated process is not modified and events are handled on demand by calling
  the appropriate ``refresh'' member function.

  \subsection{Class FollowFork}
    \declaredin{PlatFeatures.h}

\code{FollowFork} is used to configure ProcControlAPI's behavior when the
associated Process forks.

\begin{apient}
enum follow_t {None, ImmediateDetach, DisableBreakpointsDetach, Follow };
\end{apient}

\apidesc{
The follow\_t type represents the configurable behaviors under forking.
None is specified when fork tracking is not available for the current platform.
ImmediateDetach means that forked children are never attached to.
DisableBreakpointsDetach means that inherited breakpoints are removed from
forked children, and then the children are detached.
Follow is the default behavior, and it means that forked children are attached
to and remain under full control of ProcControlAPI.
}

\begin{apient}
static void setDefaultFollowFork(follow_t f);
\end{apient}

\apidesc{
This function sets the default forking behavior across all Process objects to f.
}

\begin{apient}
static follow_t getDefaultFollowFork();
\end{apient}

\apidesc{
This function returns the current default forking behavior across all Process
objects.
}

\begin{apient}
bool setFollowFork(follow_t f) const;
\end{apient}

\apidesc{
This function sets the forking behavior for the associated Process object to f.
This function returns true on success and false on failure.
}

\begin{apient}
follow_t getFollowFork() const;
\end{apient}

\apidesc{
This function returns the current forking behavior for the associated Process
object.
}

  \subsection{Class LWPTracking}
    \declaredin{PlatFeatures.h}

\code{LWPTracking} is used to configure the handling of LWP events for its
associated Process.

\begin{apient}
static void setDefaultTrackLWPs(bool b);
\end{apient}

\apidesc{
Sets the default handling mechanism for LWP events across all Process objects to
interrupt-driven (b = true) or polling-driven (b = false).
}

\begin{apient}
static bool getDefaultTrackLWPs();
\end{apient}

\apidesc{
Returns the current default handling mechanism for LWP events across all Process
objects. A return value of true indicates interrupt-driven while false
indicates polling-driven.
}

\begin{apient}
void setTrackLWPs(bool b) const;
\end{apient}

\apidesc{
Sets the LWP event handling mechanism for the associated Process object to
interrupt-driven (b = true) or polling-driven (b = false).
}

\begin{apient}
bool getTrackLWPs() const;
\end{apient}

\apidesc{
Returns the current LWP event handling mechanism for the associated Process
object. A return value of true indicates interrupt-driven while false
indicates polling-driven.
}

\begin{apient}
bool refreshLWPs();
\end{apient}

\apidesc{
Manually polls for queued LWP events to handle.
Returns true on success and false on failure.
}


  \subsection{Class LibraryTracking}
    \declaredin{PlatFeatures.h}

\code{LibraryTracking} is used to configure the handling of library events for
its associated Process.

\begin{apient}
static void setDefaultTrackLibraries(bool b);
\end{apient}

\apidesc{
Sets the default handling mechanism for library events across all Process
objects to interrupt-driven (b = true) or polling-driven (b = false).
}

\begin{apient}
static bool getDefaultTrackLibraries();
\end{apient}

\apidesc{
Returns the current default handling mechanism for library events across all
Process objects. A return value of true indicates interrupt-driven while
false indicates polling-driven.
}

\begin{apient}
bool setTrackLibraries(bool b) const;
\end{apient}

\apidesc{
Sets the library event handling mechanism for the associated Process object to
interrupt-driven (b = true) or polling-driven (b = false).
Returns true on success and false on failure.
}

\begin{apient}
bool getTrackLibraries() const;
\end{apient}

\apidesc{
Returns the current library event handling mechanism for the associated Process
object. A return value of true indicates interrupt-driven while false
indicates polling-driven.
}

\begin{apient}
bool refreshLibraries();
\end{apient}

\apidesc{
Manually polls for queued library events to handle.
Returns true on success and false on failure.
}


  \subsection{Class SignalMask}
    \declaredin{PlatFeatures.h}

\code{SignalMask} class is used to configure the signal mask for its associated
Process.

\begin{apient}
dyn_sigset_t;
\end{apient}

\apidesc{
On POSIX systems, this type is equivalent to sigset\_t.
}

\begin{apient}
static void setDefaultSigMask(dyn_sigset_t s);
\end{apient}

\apidesc{
This function sets the default signal mask across all Process objects to s.
}

\begin{apient}
static dyn_sigset_t getDefaultSigMask();
\end{apient}

\apidesc{
This function returns the current default signal mask across all Process
objects.
}

\begin{apient}
bool setSigMask(dyn_sigset_t s);
\end{apient}

\apidesc{
This function sets the signal mask for the associated Process object to s.
This function returns true on success and false on failure.
}

\begin{apient}
dyn_sigset_t getSigMask() const;
\end{apient}

\apidesc{
This function returns the current signal mask for the associated Process object.
}

  \subsection{Class ThreadTracking}
    \declaredin{PlatFeatures.h}

\code{ThreadTracking} is used to configure the handling of thread events for its
associated Process.

\begin{apient}
static void setDefaultTrackThreads(bool b);
\end{apient}

\apidesc{
Sets the default handling mechanism for thread events across all Process objects
to interrupt-driven (b = true) or polling-driven (b = false).
}

\begin{apient}
static bool getDefaultTrackThreads();
\end{apient}

\apidesc{
Returns the current default handling mechanism for thread events across all
Process objects. A return value of true indicates interrupt-driven while
false indicates polling-driven.
}

\begin{apient}
bool setTrackThreads(bool b) const;
\end{apient}

\apidesc{
Sets the thread event handling mechanism for the associated Process object to
interrupt-driven (b = true) or polling-driven (b = false).
Returns true on success and false on failure.
}

\begin{apient}
bool getTrackThreads() const;
\end{apient}

\apidesc{
Returns the current thread event handling mechanism for the associated Process
object. A return value of true indicates interrupt-driven while false
indicates polling-driven.
}

\begin{apient}
bool refreshThreads();
\end{apient}

\apidesc{
Manually polls for queued thread events to handle.
Returns true on success and false on failure.
}



\end{document}
